% !TeX spellcheck = pt_BR
\documentclass[12pt,portuguese]{article}

%----------------------------------------------------------------------
% PACOTES ESSENCIAIS E TRADUÇÃO
%----------------------------------------------------------------------
\usepackage[brazilian]{babel}     % Suporte à língua portuguesa
\usepackage[utf8]{inputenc}       % Codificação de entrada
\usepackage[T1]{fontenc}          % Codificação da fonte
\usepackage{lmodern}              % Fonte escalável

%----------------------------------------------------------------------
% PACOTES MATEMÁTICOS
%----------------------------------------------------------------------
\usepackage{amsmath, amsthm, amssymb, amstext, mathtools}
\usepackage{mathrsfs}             % Para fontes como \mathscr
\usepackage{dsfont}               % Para o indicador de conjunto \mathds{1}
\usepackage[all]{xy}              % Para diagramas comutativos
\usepackage{imakeidx}             % Criação de Lista de Símbolos / Índice

% Desigualdades sempre com símbolos inclinados
\let\leq\leqslant
\let\le\leqslant
\let\ge\geqslant
\let\geq\geqslant

%----------------------------------------------------------------------
% FORMATAÇÃO, LAYOUT E CORES
%----------------------------------------------------------------------
\usepackage{geometry}
\usepackage{varwidth}
\usepackage{graphicx}
\usepackage[pdftex]{xcolor}
\usepackage{titlesec}
\usepackage{appendix}
\usepackage{enumitem}
\usepackage{float}
\usepackage{booktabs}
\usepackage{array}
\usepackage{colortbl}
\usepackage{fontawesome5}
\usepackage{lucide-icons}
\usepackage{caption}
\usepackage{stackengine}
\captionsetup{width=0.8\textwidth}

% Definição de Cores
\definecolor{teal}{RGB}{15,118,110}
\definecolor{text}{RGB}{28,25,23}
\definecolor{muted}{RGB}{87,83,78}
\definecolor{card}{RGB}{255,255,255}
\definecolor{border}{RGB}{231,229,228}

% Definição de margens e espaçamento
\setlength{\topmargin}{-0.7in}
\setlength{\oddsidemargin}{0in}
\setlength{\textheight}{9.55in}
\setlength{\textwidth}{6.5in}
\setlength{\baselineskip}{12mm}

% Formatação dos títulos de seção
\titleformat{\section}{\Large\bfseries}{\thesection.}{0.5em}{}
\titleformat{\subsection}{\large\bfseries}{\thesubsection.}{0.5em}{}
\titleformat{\subsubsection}{\normalfont\itshape}{\thesubsubsection.}{0.6em}{}

%----------------------------------------------------------------------
% TIKZ E GRÁFICOS
%----------------------------------------------------------------------
\usepackage{tikz}
\usepackage{pgfplots}
\pgfplotsset{compat=1.17}
\usetikzlibrary{positioning,calc,fit,nfold,decorations.pathreplacing,calligraphy,tikzmark,patterns,arrows.meta,bending,hobby}

%----------------------------------------------------------------------
% REFERÊNCIAS, LINKS E METADADOS DO PDF
%----------------------------------------------------------------------
\usepackage[pdftex, colorlinks, linkcolor={red!50!black}, 
citecolor={red}, urlcolor={blue!80!black}]{hyperref}

\hypersetup{
pdftitle={Vetores Contínuos: Densidade Conjunta, Marginais e Independência},
pdfauthor={L. Cioletti},
pdfsubject={Vetores aleatórios contínuos, função de distribuição conjunta, densidade conjunta, densidades marginais e independência},
pdfkeywords={vetores aleatórios contínuos, distribuição conjunta, densidade conjunta, densidades marginais, independência, MAT-UnB}
}

% --- PACOTE E COMANDO PARA A LICENÇA CREATIVE COMMONS ---
\usepackage[type={CC},modifier={by-nc-sa},version={4.0}]{doclicense}
\newcommand\blfootnote[1]{%
\begingroup
\renewcommand\thefootnote{}\footnote{#1}%
\addtocounter{footnote}{-1}%
\endgroup
}

%----------------------------------------------------------------------
% ALIASCNT E AMBIENTES DE TEOREMA (Cleveref)
%----------------------------------------------------------------------
\usepackage{aliascnt}
\theoremstyle{definition}
\newtheorem{theorem}{Teorema}

\newaliascnt{definition}{theorem}
\newtheorem{definition}[definition]{Definição}
\aliascntresetthe{definition}

\newaliascnt{example}{theorem}
\newtheorem{example}[example]{Exemplo}
\aliascntresetthe{example}

\newaliascnt{proposition}{theorem}
\newtheorem{proposition}[proposition]{Proposição}
\aliascntresetthe{proposition}

\newaliascnt{corollary}{theorem}
\newtheorem{corollary}[corollary]{Corolário}
\aliascntresetthe{corollary}

\newaliascnt{lemma}{theorem}
\newtheorem{lemma}[lemma]{Lema}
\aliascntresetthe{lemma}

\newaliascnt{remark}{theorem}
\newtheorem{remark}[remark]{Observação}
\aliascntresetthe{remark}

\newtheoremstyle{plainexercise}
  {3pt}{3pt}{\normalfont}{}{\bfseries}{.}{.5em}
  {\thmnumber{#2}\thmnote{ #3}}
\theoremstyle{plainexercise}
\newtheorem{exercise}{Exercício}

% Versão marcada dos exercícios: a estrela fica pendurada na margem esquerda,
% como nas notas anteriores.
\newtheoremstyle{starredexercise}
  {3pt}{3pt}{\normalfont}{}{\bfseries}{.}{.5em}
  {\llap{\textasteriskcentered\,}\thmnumber{#2}\thmnote{ #3}}
\theoremstyle{starredexercise}
\newtheorem{starredexercise}[exercise]{Exercício}
\theoremstyle{definition}

% Ambientes não numerados usados no texto
\newtheorem*{property}{Propriedade}
\newtheorem*{solution}{Solução}

% Se um enumerate for o primeiro conteúdo de um teorema/exercício (cabeçalho
% ainda pendente, \@inlabel true), o \leavevmode dispara o \everypar do amsthm
% que coloca o cabeçalho na própria linha, antes da lista iniciar.
\makeatletter
\setlist[enumerate]{before={\if@inlabel\leavevmode\fi}}
\makeatother

% Nomes em português para os ambientes
\AtBeginDocument{\renewcommand{\proofname}{\normalfont\textbf{Prova}}}
\renewcommand{\qedsymbol}{$\blacksquare$}

% Cleveref (deve vir depois de hyperref e amsthm)
\usepackage[brazilian,nameinlink]{cleveref}

\crefname{theorem}{teorema}{teoremas}
\Crefname{theorem}{Teorema}{Teoremas}
\crefname{definition}{definição}{definições}
\Crefname{definition}{Definição}{Definições}
\crefname{lemma}{lema}{lemas}
\Crefname{lemma}{Lema}{Lemas}
\crefname{proposition}{proposição}{proposições}
\Crefname{proposition}{Proposição}{Proposições}
\crefname{corollary}{corolário}{corolários}
\Crefname{corollary}{Corolário}{Corolários}
\crefname{example}{exemplo}{exemplos}
\Crefname{example}{Exemplo}{Exemplos}
\crefname{exercise}{exercício}{exercícios}
\Crefname{exercise}{Exercício}{Exercícios}
\crefname{remark}{observação}{observações}
\Crefname{remark}{Observação}{Observações}
\crefname{section}{seção}{seções}
\Crefname{section}{Seção}{Seções}

%----------------------------------------------------------------------
% MACROS ADICIONAIS
%----------------------------------------------------------------------
\let\oldemptyset\emptyset
\let\emptyset\varnothing
\graphicspath{{./}{../arquivos-aux/}}

% Seno na convenção em português
\renewcommand{\sin}{\operatorname{sen}}

%----------------------------------------------------------------------
% DOCUMENTO PRINCIPAL
%----------------------------------------------------------------------
\begin{document}
% --- CABEÇALHO COM O LOGO DA UnB ---
\vspace*{-0.8cm}
\begin{figure}[H]
\includegraphics{logo-mat-unb.pdf}
\end{figure}


% --- TÍTULO, AUTOR E DATA ---
\vspace*{0.3cm}
\begin{center}
{\LARGE\bf Vetores Contínuos: Densidade Conjunta, Marginais e Independência}
\vspace{1cm}

{\large \bf L. Cioletti}
\vspace{0.3cm}

{\footnotesize Setembro de 2026}
\end{center}

% --- RESUMO (ABSTRACT) ---
\begin{abstract}
Nestas notas, damos continuidade ao estudo das variáveis aleatórias
e iniciamos o estudo dos vetores aleatórios contínuos. Apresentamos a
função de distribuição conjunta e suas propriedades, definimos a
densidade conjunta e mostramos como calcular probabilidades por meio de
integração sobre regiões do plano. Em seguida, introduzimos as densidades
marginais, obtidas integrando a densidade conjunta em relação à outra
variável, e estudamos o conceito de independência de variáveis
aleatórias absolutamente contínuas, caracterizada pela fatoração da
densidade conjunta, em quase todo ponto, como produto das densidades
marginais. Desenvolvemos três aplicações à inteligência artificial:
sorteios com probabilidades softmax por meio de relógios exponenciais,
classificação baseada na informação conjunta de duas características e
escassez de amostras locais em dimensão alta. Para preparar a primeira
aplicação, estudamos a distribuição exponencial e sua caracterização pela
propriedade de falta de memória. A apresentação dos conceitos
segue de perto as
Seções 6.1 a 6.3 do Capítulo 6 de \cite{GrimmettWelsh2014}.
\end{abstract}

% --- LICENÇA NO RODAPÉ ---
\blfootnote{%
\doclicenseImage[imagewidth=2.6em]\ 
\begingroup
\hypersetup{urlcolor=darkgray}
\href{\doclicenseURL}{%
\doclicenseName - Attribution-NonCommercial-ShareAlike 4.0 International%
}%
\endgroup
}

%----------------------------------------------------------------------
% CONTEÚDO
%----------------------------------------------------------------------

\section{Vetores aleatórios e independência}\label{sec:vetores-aleatorios}
Dadas duas variáveis aleatórias $X$ e $Y$, definidas em um espaço de
probabilidade $(\Omega, \mathcal{F}, \mathbb{P})$, é frequentemente útil
pensar nelas conjuntamente como um vetor aleatório $(X, Y)$ que assume valores
em $\mathbb{R}^2$. Quando $X$ e $Y$ são discretas, podemos estudar esse
vetor aleatório por meio da função de probabilidade conjunta de $X$ e $Y$,
mas, para variáveis aleatórias gerais, essa função pode não determinar
a distribuição. Uma ferramenta aplicável em todos esses casos é a
\emph{função de distribuição conjunta}, definida a seguir.

\begin{definition}[Função de distribuição conjunta]\label{def:funcao-distribuicao-conjunta}
  A função de distribuição conjunta do par de variáveis aleatórias $X$ e $Y$
  é a função $F_{X,Y} : \mathbb{R}^2 \to [0, 1]$ dada por
  \begin{equation}
  \label{eq:def-distribuicao-conjunta}
    F_{X,Y}(x, y) = \mathbb{P}(X \leqslant x, Y \leqslant y).
  \end{equation}
\end{definition}

A notação $\mathbb{P}(X \leqslant x, Y \leqslant y)$ representa a
probabilidade da interseção dos seguintes eventos $\{X \leqslant x\}$ e
$\{Y \leqslant y\}$.

Algumas propriedades elementares das funções de distribuição conjunta
são análogas às das funções de distribuição usuais. Por exemplo,
\begin{equation}
\label{eq:lim-dist-conjunta-menos-infinito}
  \lim_{x,y \to -\infty} F_{X,Y}(x, y) = 0
\end{equation}
e
\begin{equation}
\label{eq:lim-dist-conjunta-mais-infinito}
  \lim_{x,y \to \infty} F_{X,Y}(x, y) = 1.
\end{equation}

Para justificar o primeiro limite, basta observar que
$0 \leqslant F_{X,Y}(x, y) \leqslant F_X(x)$ e que
$F_X(x) \to 0$ quando $x \to -\infty$. Para o segundo, usamos a
continuidade da probabilidade em sequências crescentes: os eventos
$\{X \leqslant n, Y \leqslant n\}$ crescem para $\Omega$ quando
$n \to \infty$, pois $X$ e $Y$ assumem valores reais. Além disso,
$F_{X,Y}(x, y) \geqslant F_{X,Y}(n, n)$ sempre que $x, y \geqslant n$.

Da mesma forma, a inclusão dos eventos mostra que $F_{X,Y}$ é
não decrescente em cada variável, no sentido de que
\begin{equation}
\label{eq:monotonicidade-distribuicao-conjunta}
  F_{X,Y}(x_1, y_1) \leqslant F_{X,Y}(x_2, y_2)
  \quad \text{se} \quad x_1 \leqslant x_2 \text{ e } y_1 \leqslant y_2.
\end{equation}

A função de distribuição conjunta $F_{X,Y}$ contém, em geral, mais informação
do que as duas funções de distribuição usuais $F_X$ e $F_Y$, uma vez que
ela descreve como $X$ e $Y$ se comportam \emph{conjuntamente}. Em
particular, as funções de distribuição de $X$ e $Y$ podem ser obtidas a
partir de sua função de distribuição conjunta de maneira direta.
Para obter $F_X(x)$, fixamos $x$ e deixamos o limiar de $Y$ tender ao
infinito. Os eventos $\{X \leqslant x, Y \leqslant n\}$ crescem até
$\{X \leqslant x\}$ quando $n \to \infty$. Pela continuidade da
probabilidade em sequências crescentes, e analogamente para $Y$, temos
\begin{equation}
\label{eq:distribuicoes-marginais-limites}
  F_X(x) = \lim_{y \to \infty} F_{X,Y}(x, y),
  \qquad
  F_Y(y) = \lim_{x \to \infty} F_{X,Y}(x, y).
\end{equation}

Essas funções de distribuição são chamadas de funções de distribuição
\emph{marginais} da função de distribuição conjunta $F_{X,Y}$.

A noção de independência das variáveis aleatórias $X$ e $Y$ surge
naturalmente dessa discussão.

\begin{definition}[Variáveis aleatórias independentes]\label{def:variaveis-aleatorias-independentes}
  Dizemos que $X$ e $Y$ são \textbf{independentes} se, para quaisquer $x, y \in
  \mathbb{R}$, os eventos $\{X \leqslant x\}$ e $\{Y \leqslant y\}$ forem
  independentes.
\end{definition}

Isto é, $X$ e $Y$ são independentes se, e somente se,
\[
  \mathbb{P}(X \leqslant x, Y \leqslant y)
  = \mathbb{P}(X \leqslant x)\,\mathbb{P}(Y \leqslant y)
  \quad \text{para todos } x, y \in \mathbb{R},
\]
o que equivale a dizer que a sua função de distribuição conjunta se fatora
como o produto das duas funções de distribuição marginais:
\begin{equation}
\label{eq:fatoracao-distribuicao-conjunta}
  F_{X,Y}(x, y) = F_X(x)\,F_Y(y)
  \quad \text{para todos } x, y \in \mathbb{R}.
\end{equation}

No caso de variáveis aleatórias \emph{discretas}, essa definição equivale
à fatoração da função de probabilidade conjunta como produto das funções
de probabilidade marginais. A verificação fica como exercício.
Variáveis aleatórias que não são independentes são chamadas de
\emph{dependentes}.

Estudamos famílias finitas de variáveis aleatórias da mesma maneira.
Se $X_1, X_2, \ldots, X_n$ são variáveis aleatórias em
$(\Omega, \mathcal{F}, \mathbb{P})$, escrevemos
$\mathbf{X} = (X_1, X_2, \ldots, X_n)$. A sua \emph{função de
distribuição conjunta} é a função
$F_{\mathbf{X}} : \mathbb{R}^n \to [0, 1]$ dada por
\begin{equation}
\label{eq:distribuicao-conjunta-n-dimensoes}
  F_{\mathbf{X}}(\mathbf{x})
  = \mathbb{P}(X_1 \leqslant x_1, X_2 \leqslant x_2, \ldots, X_n \leqslant x_n),
  \qquad \forall \mathbf{x} = (x_1, x_2, \ldots, x_n) \in \mathbb{R}^n.
\end{equation}
Vamos dizer que uma coleção finita de variáveis aleatórias
$X_1, X_2, \ldots, X_n$ são \emph{independentes} se
\[
  \mathbb{P}(X_1 \leqslant x_1, \ldots, X_n \leqslant x_n)
  = \mathbb{P}(X_1 \leqslant x_1)\cdots\mathbb{P}(X_n \leqslant x_n)
  \qquad \forall \mathbf{x} \in \mathbb{R}^n,
\]
ou, equivalentemente, se
\begin{equation}
\label{eq:fatoracao-distribuicao-n-dimensoes}
  F_{\mathbf{X}}(\mathbf{x})
  = F_{X_1}(x_1)\cdots F_{X_n}(x_n)
  \qquad \forall\mathbf{x} \in \mathbb{R}^n.
\end{equation}

\begin{example}[Função de distribuição conjunta discreta]\label{ex:distribuicao-conjunta-discreta}
Suponha que $X$ e $Y$ sejam variáveis aleatórias em um espaço de probabilidade,
cada uma assumindo valores no conjunto dos inteiros
$\mathbb{Z} = \{\ldots, -1, 0, 1, \ldots\}$,
com função de probabilidade conjunta
\[
  \mathbb{P}(X = i, Y = j) = p(i, j)
  \quad\text{para } i, j \in \mathbb{Z}.
\]
Somando as probabilidades dos eventos $\{X = i, Y = j\}$,
que são dois a dois disjuntos, obtemos
\[
  F_{X,Y}(x, y)
  = \sum_{\substack{i,j\in\mathbb{Z} \\ i\leqslant x,\ j\leqslant y}} p(i, j)
  \quad\text{para todo } (x, y) \in \mathbb{R}^2.
\]
\end{example}

\begin{example}[Distribuição conjunta e independência de exponenciais]\label{ex:distribuicao-conjunta-exponenciais}
Suponha que $X$ e $Y$ sejam variáveis aleatórias com função de distribuição conjunta
\[
  F_{X,Y}(x, y) =
  \begin{cases}
    1 - e^{-x} - e^{-y} + e^{-x-y} & \text{se } x, y \geqslant 0,
    \\[0.2cm]
    0 & \text{caso contrário.}
  \end{cases}
\]
A função de distribuição marginal de $X$ é
\[
  F_X(x) = \lim_{y \to \infty} F_{X,Y}(x, y) =
  \begin{cases}
    1 - e^{-x} & \text{se } x \geqslant 0,
    \\[0.2cm]
    0 & \text{caso contrário,}
  \end{cases}
\]
de modo que $X$ tem distribuição exponencial com parâmetro $1$. Um cálculo
análogo mostra que $Y$ também possui essa distribuição.
Para $x, y \geqslant 0$, podemos fatorar a expressão inicial:
\[
  F_{X,Y}(x, y) = (1 - e^{-x})(1 - e^{-y}) = F_X(x)\,F_Y(y).
\]
Se $x < 0$ ou $y < 0$, tanto $F_{X,Y}(x, y)$ quanto
$F_X(x)\,F_Y(y)$ são nulos. Portanto, a fatoração vale para todos
$x, y \in \mathbb{R}$, e $X$ e $Y$ são independentes.
\end{example}

\begin{exercise}\label{exercise:independencia-probabilidades-cauda}
Mostre que duas variáveis aleatórias $X$ e $Y$ são independentes se, e somente se,
\[
  \mathbb{P}(X > x, Y > y) = \mathbb{P}(X > x)\,\mathbb{P}(Y > y)
  \quad\text{para todos } x, y \in \mathbb{R}.
\]
\end{exercise}

\begin{exercise}\label{exercise:probabilidade-retangulo-fdc}
Seja $(X, Y)$ um par de variáveis aleatórias com função de distribuição
conjunta $F$.
Prove que
\[
  \mathbb{P}(a < X \leqslant b, \, c < Y \leqslant d)
  = F(b, d) + F(a, c) - F(a, d) - F(b, c)
\]
para quaisquer $a, b, c, d \in \mathbb{R}$ tais que $a < b$ e $c < d$.
\end{exercise}

\begin{exercise}\label{exercise:independencia-retangulos}
Prove que duas variáveis aleatórias $X$ e $Y$ são independentes se, e somente se,
\[
  \mathbb{P}(a < X \leqslant b, \, c < Y \leqslant d)
  = \mathbb{P}(a < X \leqslant b)\,\mathbb{P}(c < Y \leqslant d)
\]
para todos $a, b, c, d \in \mathbb{R}$ satisfazendo $a < b$ e $c < d$.
\end{exercise}

\section{Funções de densidade conjunta}\label{sec:densidade-conjunta}
Uma densidade conjunta permitirá calcular a probabilidade de o vetor
aleatório pertencer a uma região do plano. Antes de defini-la, vamos
especificar quais regiões serão mensuráveis, estendendo ao plano a noção
de conjunto boreliano já apresentada na reta.

Dados $a \in \mathbb{R}^2$ e $r > 0$, o disco aberto de centro $a$ e raio
$r$ é o conjunto
\[
  D(a,r) = \{x \in \mathbb{R}^2 : \|x-a\| < r\},
\]
isto é, o conjunto dos pontos cuja distância até $a$ é menor que $r$.
Aqui, $\|\cdot\|$ denota a norma euclidiana.
\begin{figure}[H]
\centering
\begin{tikzpicture}[scale=0.85, >=Stealth]
  \draw[->] (-0.35,0) -- (4.5,0) node[below right] {$x_1$};
  \draw[->] (0,-0.35) -- (0,3.65) node[above left] {$x_2$};
  \fill[teal!10] (2.1,1.65) circle (1.25);
  \draw[teal, dashed, line width=0.9pt] (2.1,1.65) circle (1.25);
  \fill (2.1,1.65) circle (1.5pt) node[below left] {$a$};
  \draw[teal!80!black, line width=0.8pt] (2.1,1.65) -- (2.984,2.534);
  \node[teal!80!black, above left] at (2.56,2.11) {$r$};
  \node[teal!80!black] at (2.05,3.3) {$D(a,r)$};
\end{tikzpicture}
\caption{Disco aberto de centro $a$ e raio $r$ em $\mathbb{R}^2$.}
\label{fig:disco-aberto}
\end{figure}
Um conjunto $U \subseteq \mathbb{R}^2$ é \emph{aberto} se, para todo ponto
$a \in U$, existe algum raio $r > 0$ tal que $D(a,r) \subseteq U$.
Denotemos por $\mathscr{C}$ a coleção de todos os discos abertos de
$\mathbb{R}^2$, com centros em $\mathbb{R}^2$ e raios positivos.
A $\sigma$-álgebra de Borel de
$\mathbb{R}^2$, denotada por $\mathscr{B}(\mathbb{R}^2)$, é a menor
$\sigma$-álgebra de subconjuntos de $\mathbb{R}^2$ que contém $\mathscr{C}$. Em
outras palavras, é a interseção de todas as $\sigma$-álgebras de subconjuntos
de $\mathbb{R}^2$ que contêm todos os discos abertos. Seus elementos são
chamados de \emph{conjuntos borelianos}, ou simplesmente \emph{borelianos}.

Todo aberto do plano é uma união enumerável de discos abertos com centros
em $\mathbb{Q}^2$ e raios racionais positivos. Por isso, a classe dos
borelianos contém todos os conjuntos abertos e fechados do plano. Em
particular, inclui retângulos, discos e regiões delimitadas por curvas
fechadas simples (curvas de Jordan), entre outros conjuntos que aparecem
nas aplicações do Cálculo.

Diremos que uma função $f:\mathbb{R}^2\to\mathbb{R}$ é
\emph{Borel-mensurável} se, para todo $t \in \mathbb{R}$, o conjunto
\[
  \{(x,y) \in \mathbb{R}^2 : f(x,y) \leqslant t\}
\]
pertencer a $\mathscr{B}(\mathbb{R}^2)$. Como os intervalos $(-\infty,t]$
geram $\mathscr{B}(\mathbb{R})$, essa condição equivale a exigir que a
pré-imagem $f^{-1}(C) = \{(x,y) \in \mathbb{R}^2 : f(x,y) \in C\}$
pertença a $\mathscr{B}(\mathbb{R}^2)$ para todo
$C \in \mathscr{B}(\mathbb{R})$. Essa condição é análoga à usada na
definição de variável aleatória: as pré-imagens dos borelianos da reta
devem ser conjuntos mensuráveis no domínio. Na reta, usamos a mesma
definição, tomando $\mathscr{B}(\mathbb{R})$ também no domínio.
Em particular, toda função
contínua é Borel-mensurável, pois a pré-imagem de um aberto por uma função
contínua é um conjunto aberto.

Com essas noções, podemos passar das densidades na reta às densidades no plano.
Recordamos que uma variável aleatória $X$ é dita \emph{absolutamente contínua}
se sua função de distribuição admitir uma representação integral da forma
\[
  F_X(x) = \mathbb{P}(X \leqslant x) = \int_{-\infty}^{x} f(u)\,du,
  \qquad \text{para todo } x \in \mathbb{R},
\]
para alguma função Borel mensurável e não negativa $f$.

\begin{definition}[Variáveis aleatórias conjuntamente absolutamente contínuas]\label{def:va-conjuntamente-continuas}
  Dizemos que o par $(X, Y)$ de variáveis aleatórias em
  $(\Omega, \mathcal{F}, \mathbb{P})$ é
  \emph{(conjuntamente) absolutamente contínuo} se sua função de distribuição
  conjunta puder ser expressa na forma\footnote{As integrais são tomadas no
  sentido de Lebesgue. Nos exemplos do curso, elas podem ser calculadas pelas
  técnicas usuais do Cálculo. Como o integrando é mensurável e não negativo,
  o Teorema de Tonelli permite trocar a ordem de integração. Quando sua
  integral é finita, também se aplica o Teorema de Fubini.}
  \begin{equation}\label{eq:distr-conjunta-integral}
    F_{X,Y}(x, y) = \mathbb{P}(X \leqslant x, Y \leqslant y)
    = \int_{-\infty}^{x} \int_{-\infty}^{y} f(u, v)\,dv\,du
  \end{equation}
  para quaisquer $x, y \in \mathbb{R}$ e alguma função Borel-mensurável
  $f : \mathbb{R}^2 \to [0, \infty)$. Quando essa condição é satisfeita,
  dizemos que $X$ e $Y$ possuem \emph{função densidade (de probabilidade)
  conjunta} $f$, e costumamos denotar essa função por $f_{X,Y}$.
\end{definition}

A medida de Lebesgue no plano formaliza a ideia de área, assim como na reta
formaliza o comprimento. Um conjunto tem \emph{área zero} se, para todo
$\varepsilon > 0$, puder ser coberto por uma união enumerável de retângulos
cuja soma das áreas seja menor que $\varepsilon$. Por exemplo, todo segmento
de reta tem área zero. Dizemos que uma propriedade vale
\emph{quase em todo ponto} do plano quando ela pode falhar apenas em um
conjunto de área zero.

Assim como no caso unidimensional, duas densidades representam a mesma
distribuição conjunta se, e somente se, coincidem quase em todo ponto.
Em particular, a densidade não é única: alterar seus valores em um conjunto
de área zero, mantendo a Borel-mensurabilidade e a não negatividade, não altera
a representação integral de $F_{X,Y}$.

Sob hipóteses de regularidade, podemos obter uma densidade por derivação.
Uma condição suficiente é que o par $(X, Y)$ seja conjuntamente absolutamente
contínuo e que $F_{X,Y}$ tenha derivadas parciais de segunda ordem contínuas
em um aberto $G \subseteq \mathbb{R}^2$ cujo complemento tenha área zero.
Nesse caso, uma densidade conjunta é dada por
\begin{equation}\label{eq:densidade-derivada-mista}
  f_{X,Y}(x, y) =
  \begin{cases}
    \dfrac{\partial^2}{\partial x \partial y} F_{X,Y}(x, y)
    & \text{se } (x, y) \in G,
    \\[0.3cm]
    0 & \text{caso contrário}.
  \end{cases}
\end{equation}
Para entender essa fórmula, observe que a integral da derivada mista sobre
um retângulo cujo fecho esteja contido em $G$ é, pelo Teorema Fundamental
do Cálculo, a diferença de quatro valores de $F_{X,Y}$ que dá a probabilidade
desse retângulo (\Cref{exercise:probabilidade-retangulo-fdc}). Isso identifica a derivada
mista com uma densidade quase em todo ponto de $G$. A escolha do valor zero
fora de $G$ não altera as integrais, pois esse complemento tem área zero.
Nos exemplos em que essa hipótese estiver satisfeita, usaremos
\eqref{eq:densidade-derivada-mista} para encontrar uma densidade conjunta.

As propriedades elementares da função densidade conjunta $f_{X,Y}$ são a
não negatividade, exigida na definição,
\begin{equation}\label{eq:densidade-nao-negativa}
  f_{X,Y}(x, y) \geqslant 0 \quad \text{para quaisquer } x, y \in \mathbb{R},
\end{equation}
e a integral unitária,
\begin{equation}\label{eq:densidade-integral-unitaria}
  \int_{-\infty}^{\infty} \int_{-\infty}^{\infty} f_{X,Y}(x, y)\,dx\,dy = 1.
\end{equation}
Esta última igualdade é obtida tomando o limite quando $x$ e $y$ tendem a $+\infty$ na igualdade
\eqref{eq:distr-conjunta-integral}, já que os eventos
$\{X \leqslant n, Y \leqslant n\}$ crescem para $\Omega$ e, pela não
negatividade de $f$, as integrais correspondentes crescem para a integral sobre
todo o plano.

Mais uma vez, observamos uma analogia entre funções densidade conjuntas e
funções de probabilidade conjuntas. De modo intuitivo, podemos dizer que, para
um ponto $(x, y)$ em que $f_{X,Y}$ seja contínua e para acréscimos positivos
pequenos $\delta x$ e $\delta y$, a probabilidade de o vetor aleatório $(X, Y)$
pertencer ao pequeno retângulo com vértice inferior esquerdo em $(x, y)$ e
lados de comprimentos $\delta x$ e $\delta y$ é dada aproximadamente por
\begin{equation}\label{eq:prob-retangulo-infinitesimal}
  \mathbb{P}\bigl(x < X \leqslant x + \delta x,\  y < Y \leqslant y + \delta y\bigr)
  \approx f_{X,Y}(x, y)\,\delta x\,\delta y.
\end{equation}
De fato, pela continuidade, a densidade varia pouco nesse retângulo.
O erro da aproximação, dividido pela área $\delta x\,\delta y$, tende a zero
quando $\delta x$ e $\delta y$ tendem a zero por valores positivos.
Essa interpretação sugere que a probabilidade de uma região deve ser obtida
integrando a densidade sobre ela, como veremos a seguir.

\begin{theorem}[Probabilidade de regiões borelianas]\label{thm:probabilidade-subconjuntos-regulares}
  Se $A$ é um boreliano de $\mathbb{R}^2$ e $X$ e $Y$ são
  variáveis aleatórias conjuntamente absolutamente contínuas com função
  densidade conjunta $f_{X,Y}$, então
  \begin{equation}\label{eq:probabilidade-conjunta-regiao}
    \mathbb{P}\bigl((X, Y) \in A\bigr) = \iint_A f_{X,Y}(x, y)\,dx\,dy.
  \end{equation}
\end{theorem}
O evento $\{(X, Y) \in A\}$ pertence a $\mathcal{F}$. Para justificar
essa afirmação, observamos que a pré-imagem de um retângulo
$(a,b]\times(c,d]$ é a interseção dos eventos
$\{a < X \leqslant b\}$ e $\{c < Y \leqslant d\}$ e, portanto, pertence
a $\mathcal{F}$. A coleção dos conjuntos do plano cujas pré-imagens
pertencem a $\mathcal{F}$ é uma $\sigma$-álgebra. Ela contém todos os
discos abertos, pois cada um deles é uma união enumerável desses
retângulos com extremos racionais. Logo, contém
$\mathscr{B}(\mathbb{R}^2)$, o que justifica a afirmação para todo boreliano $A$.

\begin{figure}[H]
\centering
\begin{tikzpicture}[scale=1.75, >=Stealth,
  regiao/.style={use Hobby shortcut, closed, draw=teal,
    fill=teal!10, line width=0.8pt},
]
  \draw[->] (-0.4,0) -- (5.25,0) node[below] {$x$};
  \draw[->] (0,-0.4) -- (0,4.05) node[left] {$y$};
  % Mesmo contorno suave da região na figura de pré-imagem das notas anteriores.
  \draw[regiao]
    (0.52,2.50) .. (1.62,0.66) .. (2.48,1.14) .. (3.60,0.54)
    .. (4.41,1.70) .. (4.61,2.66) .. (4.10,3.19)
    .. (2.99,3.37) .. (1.87,3.51);
  \node[teal!80!black] at (3.72,2.70) {$A$};

  \draw[dashed, gray!70] (2.05,0) -- (2.05,1.60);
  \draw[dashed, gray!70] (2.75,0) -- (2.75,2.00);
  \draw[dashed, gray!70] (0,1.60) -- (2.05,1.60);
  \draw[dashed, gray!70] (0,2.00) -- (2.75,2.00);
  \draw[fill=teal!40, draw=teal, line width=0.6pt]
    (2.05,1.60) rectangle (2.75,2.00);
  \node[below] at (2.05,0) {$x$};
  \node[below] at (2.95,0) {$x+\delta x$};
  \node[left] at (0,1.55) {$y$};
  \node[left] at (0,2.05) {$y+\delta y$};
  \node[teal!80!black] (area) at (3.50,2.10) {$\delta x\,\delta y$};
  \draw[->, teal!80!black, line width=0.45pt]
    (area.west) to[out=190,in=30] (2.50,1.82);
\end{tikzpicture}
\caption{Quando a densidade varia pouco no retângulo destacado em $A$,
a probabilidade de que $(X, Y)$ pertença a ele é aproximadamente
$f_{X,Y}(x, y)\,\delta x\,\delta y$.}
\label{fig:densidade-conjunta-interpretacao}
\end{figure}

O \Cref{thm:probabilidade-subconjuntos-regulares} depende de resultados da
teoria da integração, e omitimos sua demonstração.

Podemos entender intuitivamente o resultado para uma região limitada $A$
como a da \Cref{fig:densidade-conjunta-interpretacao}. Suponha que sua fronteira
tenha área zero e que a densidade seja contínua em uma vizinhança de seu
fecho. A ideia é aproximar $A$ por uniões de pequenos retângulos dois a dois
disjuntos. Em cada retângulo, a probabilidade é aproximadamente a densidade
em um de seus pontos multiplicada pela área. Somando essas pequenas
probabilidades e refinando a divisão, obtemos somas que convergem para a
integral da densidade sobre $A$.

\begin{example}[Distribuição uniforme em um retângulo]\label{ex:distribuicao-uniforme-retangulo}
  Uma função Borel-mensurável $f : \mathbb{R}^2 \to \mathbb{R}$
  é a função densidade conjunta de algum par de variáveis aleatórias,
  em algum espaço de probabilidade, se, e somente se, satisfaz as condições
  \[
    f(x, y) \geqslant 0 \quad \text{para quaisquer } x, y \in \mathbb{R}
    \quad \text{e} \quad
    \int_{-\infty}^{\infty} \int_{-\infty}^{\infty} f(x, y)\,dx\,dy = 1.
  \]
  Para verificar a suficiência, podemos tomar $\Omega = \mathbb{R}^2$ e
  $\mathcal{F} = \mathscr{B}(\mathbb{R}^2)$ e definir
  \[
    \mathbb{P}(A) = \iint_A f(x,y)\,dx\,dy
    \quad \text{para todo } A \in \mathcal{F}.
  \]
  Usamos então as variáveis aleatórias $X(x,y) = x$ e $Y(x,y) = y$.
  A mensurabilidade, a não negatividade e a
  integral unitária garantem que essa construção define uma medida de
  probabilidade. A caracterização é análoga à do caso unidimensional.

  Fixe agora $a > 0$ e $b > 0$. A função
  \[
    f(x, y) =
    \begin{cases}
      \dfrac{1}{ab}, & \text{se } 0 < x < a \text{ e } 0 < y < b, \\[0.3cm]
      0, & \text{caso contrário},
    \end{cases}
  \]
  é mensurável, não negativa e satisfaz
  \[
    \int_0^a \int_0^b \frac{1}{ab}\,dy\,dx = 1.
  \]
  Portanto, é uma função densidade conjunta. Se $X$ e $Y$ têm densidade
  conjunta $f$, dizemos que o vetor $(X, Y)$ é uniformemente distribuído no retângulo
  $B = (0, a) \times (0, b)$. Para qualquer boreliano $A$ do plano,
  \begin{align*}
    \mathbb{P}\bigl((X, Y) \in A\bigr)
    = \iint_A f(x, y)\,dx\,dy
    = \iint_{A \cap B} \frac{1}{ab}\,dx\,dy
    = \frac{\text{área}(A \cap B)}{\text{área}(B)}.
  \end{align*}
\end{example}

\begin{figure}[H]
\centering
\begin{tikzpicture}[
  x={(1.15cm,0cm)}, y={(0.65cm,0.38cm)}, z={(0cm,1.05cm)},
  >=Stealth, line join=round, line cap=round,
  eixo/.style={->, line width=0.5pt},
  face/.style={draw=teal!70!black, line width=0.8pt},
]
  % Eixos coordenados
  \draw[eixo] (-0.45,0,0) -- (4.65,0,0) node[below right] {$x$};
  \draw[eixo] (0,-0.45,0) -- (0,4.25,0) node[above left] {$y$};
  \draw[eixo] (0,0,0) -- (0,0,3.85) node[above=10pt,left] {$f(x,y)$};

  % Prisma sob a densidade constante em B = (0,a) x (0,b)
  \fill[teal!18] (0,0,2) -- (4,0,2) -- (4,3,2) -- (0,3,2) -- cycle;
  \draw[face] (0,0,0) -- (4,0,0) -- (4,3,0) -- (0,3,0) -- cycle;
  \draw[face] (0,0,2) -- (4,0,2) -- (4,3,2) -- (0,3,2) -- cycle;
  \draw[face] (0,0,0) -- (0,0,2);
  \draw[face] (4,0,0) -- (4,0,2);
  \draw[face] (4,3,0) -- (4,3,2);
  \draw[face, densely dashed] (0,3,0) -- (0,3,2);

  % Marcas dos comprimentos e da altura
  \fill (4,0,0) circle (2pt) node[below,xshift=-5pt] {$a$};
  \fill (0,3,0) circle (2pt)
    node[left, above=-5pt, xshift=-8pt] {$b$};
  \fill (0,0,2) circle (2pt);
  \draw[densely dashed] (0,0,2) -- (0,0,0);
  \node[left] at (0,0,2) {$\dfrac{1}{ab}$};
\end{tikzpicture}
\caption{A densidade uniforme tem altura $1/(ab)$ sobre o retângulo
$B=(0,a)\times(0,b)\subseteq \mathbb{R}^2$. O prisma representado tem volume $1$.}
\label{fig:densidade-uniforme-retangulo}
\end{figure}

\begin{exercise}\label{exercise:constante-e-distribuicao-conjunta}
  As variáveis aleatórias $X$ e $Y$ têm função densidade conjunta dada por
  \[
    f(x, y) =
    \begin{cases}
      c\bigl(x^2 + \frac{1}{2}xy\bigr), & \text{se } 0 < x < 1 \text{ e } 0 < y < 2, \\[0.3cm]
      0, & \text{caso contrário}.
    \end{cases}
  \]
  Determine o valor da constante $c$ e a função de distribuição conjunta
  de $X$ e $Y$.
\end{exercise}

\begin{exercise}\label{exercise:probabilidades-densidade-exponencial}
  As variáveis aleatórias $X$ e $Y$ têm função densidade conjunta dada por
  \[
    f(x, y) =
    \begin{cases}
      e^{-x-y}, & \text{se } x > 0 \text{ e } y > 0, \\[0.3cm]
      0, & \text{caso contrário}.
    \end{cases}
  \]
  Calcule $\mathbb{P}(X + Y \leqslant 1)$ e $\mathbb{P}(X > Y)$.
\end{exercise}

\section{Densidades marginais e independência}\label{sec:densidades-marginais}
Sempre que o par $(X, Y)$ possui função densidade conjunta $f_{X,Y}$,
podemos obter as funções densidade de $X$ e de $Y$ integrando em relação
à outra coordenada. Para encontrar uma densidade de $X$, observamos que
\[
  \mathbb{P}(X \leqslant x)
  = \int_{-\infty}^{x}
  \left[\int_{-\infty}^{\infty} f_{X,Y}(u, v)\,dv\right]\,du.
\]
A integração iterada é válida porque a densidade conjunta é não negativa.
Além disso, a função entre colchetes é mensurável, não negativa e tem
integral $1$ sobre a reta. Assim, uma densidade de $X$ é dada por
\begin{equation}\label{eq:densidade-marginal-x}
  f_X(x) = \int_{-\infty}^{\infty} f_{X,Y}(x, v)\,dv
  \quad \text{para quase todo ponto } x \in \mathbb{R}.
\end{equation}
De modo análogo,
\begin{equation}\label{eq:densidade-marginal-y}
  f_Y(y) = \int_{-\infty}^{\infty} f_{X,Y}(u, y)\,du
  \quad \text{para quase todo ponto } y \in \mathbb{R}.
\end{equation}
Como essas funções têm integral $1$, as integrais acima são finitas em
quase todo ponto. Se uma delas for infinita em um ponto excepcional,
fixamos a densidade marginal igual a zero nesse ponto. O conjunto de tais
pontos tem comprimento nulo, de modo que essa modificação não altera as
probabilidades.

Essas funções densidade são chamadas de funções densidade \emph{marginais}
de $X$ e de $Y$, pois são obtidas ao ``projetar'' o vetor aleatório
$(X, Y)$ sobre os dois eixos coordenados do plano.


\bigskip

Recorde, pela \Cref{def:variaveis-aleatorias-independentes},
que $X$ e $Y$ são \emph{independentes} se suas funções de distribuição
satisfazem
\begin{equation}\label{eq:independencia-distribuicao-conjunta}
  F_{X,Y}(x, y) = F_X(x)F_Y(y) \qquad \forall (x,y) \in \mathbb{R}^2.
\end{equation}

Suponha agora que $X$ e $Y$ sejam conjuntamente absolutamente contínuas.
O produto $f_X(x)f_Y(y)$ é uma função mensurável não negativa e tem
integral $1$ sobre o plano. Sua função de distribuição é
\[
  \int_{-\infty}^{x}\int_{-\infty}^{y}
  f_X(u)f_Y(v)\,dv\,du = F_X(x)F_Y(y).
\]
Se $X$ e $Y$ são independentes, essa função de distribuição coincide com
$F_{X,Y}$. Portanto, $f_{X,Y}$ e o produto das densidades marginais são
densidades da mesma distribuição e coincidem quase em todo ponto:
\begin{equation}\label{eq:independencia-densidade-conjunta}
  f_{X,Y}(x, y) = f_X(x)f_Y(y)
  \quad \text{para quase todo ponto } (x,y) \in \mathbb{R}^2.
\end{equation}
Aqui, ``quase todo ponto'' permite que a igualdade falhe em um conjunto
de área nula, pois os valores de uma densidade nesse conjunto não alteram
as probabilidades. Reciprocamente, se
integrarmos ambos os lados de \eqref{eq:independencia-densidade-conjunta}
com $x$ variando em $(-\infty, u]$ e $y$ variando em $(-\infty, v]$,
obtemos \eqref{eq:independencia-distribuicao-conjunta}. Assim, variáveis
aleatórias conjuntamente absolutamente contínuas são independentes se,
e somente se, sua função densidade conjunta se fatora, quase em todo ponto,
como o produto das duas funções densidade marginais.
Esse resultado é análogo ao caso de variáveis
aleatórias discretas, estudado anteriormente. Do mesmo modo que no caso
discreto, há um resultado mais geral.

\begin{theorem}[Critério de fatoração para independência]\label{thm:criterio-fatoracao-independencia}
Sejam $X$ e $Y$ variáveis aleatórias conjuntamente absolutamente contínuas.
Elas são independentes se, e somente se, existem funções Borel-mensuráveis
não negativas $g,h:\mathbb{R}\to[0,\infty)$ tais que
\[
  f_{X,Y}(x, y) = g(x)h(y)
  \quad \text{para quase todo ponto } (x,y) \in \mathbb{R}^2.
\]
Isto é, a densidade conjunta se fatora como o produto de uma função da
primeira coordenada por uma função da segunda.
\end{theorem}

Não apresentamos a demonstração aqui, mas sugerimos ao leitor que adapte
a demonstração correspondente dada para variáveis discretas, substituindo os
somatórios por integrais. Para essa adaptação, verifique primeiro que a
integral unitária da densidade conjunta implica que as integrais de $g$ e
de $h$ são finitas e estritamente positivas. Isso permite normalizar os
fatores, dividindo cada um por sua integral, para obter as densidades marginais.

Não vamos nos estender longamente sobre o caso em que há três ou mais
variáveis aleatórias. De modo geral, todas as ideias desenvolvidas até
aqui possuem análogos em mais de duas dimensões. Em $\mathbb{R}^3$, a
$\sigma$-álgebra de Borel $\mathscr{B}(\mathbb{R}^3)$ é gerada pelas bolas
abertas $\{x \in \mathbb{R}^3 : \|x-a\| < r\}$, com
$a \in \mathbb{R}^3$ e $r > 0$. A Borel-mensurabilidade de funções reais
é definida pelo mesmo critério de pré-imagens, usando
$\mathscr{B}(\mathbb{R}^3)$ no domínio. Volume nulo e propriedades quase
em todo ponto são definidos como no plano, substituindo os retângulos
e suas áreas por paralelepípedos e seus volumes.

Por exemplo, dizemos que
três variáveis aleatórias $X$, $Y$ e $Z$ são
\emph{conjuntamente absolutamente contínuas} se
\[
  \begin{aligned}
    F_{X,Y,Z}(x,y,z)
    &= \mathbb{P}(X \leqslant x, Y \leqslant y, Z \leqslant z)
    \\[0.3cm]
    &= \int_{-\infty}^{x}\int_{-\infty}^{y}\int_{-\infty}^{z}
    f(u, v, w)\,dw\,dv\,du
  \end{aligned}
\]
para todo $(x,y,z) \in \mathbb{R}^3$ e alguma função Borel-mensurável
não negativa $f:\mathbb{R}^3\to[0,\infty)$ cuja integral sobre
$\mathbb{R}^3$ seja $1$. Como no caso bidimensional, quando o vetor
$(X,Y,Z)$ é conjuntamente absolutamente contínuo e $F_{X,Y,Z}$
tem derivadas parciais contínuas até a ordem três em um aberto
$H \subseteq \mathbb{R}^3$ cujo complemento tenha volume nulo,
podemos obter uma densidade pela fórmula
\[
  f(x, y, z) =
  \begin{cases}
    \dfrac{\partial^3 F_{X,Y,Z}}{
    \partial x\,\partial y\,\partial z}(x,y,z),
    & \text{se } (x,y,z) \in H, \\[0.3cm]
    0, & \text{caso contrário}.
  \end{cases}
\]
Para qualquer vetor conjuntamente absolutamente contínuo $(X,Y,Z)$,
as três variáveis são \emph{independentes} se, e somente se, a densidade
conjunta se fatora como o produto das funções densidade marginais:
\[
  f(x, y, z) = f_X(x)f_Y(y)f_Z(z)
  \quad \text{para quase todo ponto } (x,y,z) \in \mathbb{R}^3.
\]
Nesse caso, o conjunto excepcional tem volume nulo.

\begin{example}[Exponenciais independentes]\label{ex:exponenciais-independentes}
  Suponha que $X$ e $Y$ tenham função densidade conjunta dada por
  \[
    f(x, y) =
    \begin{cases}
      e^{-x-y}, & \text{se } x > 0 \text{ e } y > 0, \\[0.2cm]
      0, & \text{caso contrário}.
    \end{cases}
  \]
  Então
  \[
    \begin{aligned}
      f_X(x)
      &= \int_{-\infty}^{\infty} f(x, y)\,dy
      \\[0.3cm]
      &= \begin{cases}
        \displaystyle\int_0^{\infty} e^{-x-y}\,dy,
        & \text{se } x > 0, \\[0.2cm]
        0, & \text{caso contrário},
      \end{cases}
      \\[0.3cm]
      &= \begin{cases}
        e^{-x}, & \text{se } x > 0, \\[0.2cm]
        0, & \text{caso contrário},
      \end{cases}
    \end{aligned}
  \]
  o que mostra que $X$ tem distribuição exponencial com parâmetro $1$. A
  variável aleatória $Y$ também possui essa distribuição. Com essas
  escolhas das densidades marginais, temos
  \[
    f_{X,Y}(x, y) = f_X(x)f_Y(y) \quad \text{para todo } x, y \in \mathbb{R},
  \]
  de modo que $X$ e $Y$ são independentes.
\end{example}

\begin{example}[Constante de normalização e dependência]\label{ex:constante-normalizacao-dependencia}
  Suponha que $X$ e $Y$ tenham função densidade conjunta dada por
  \begin{equation}\label{eq:densidade-conjunta-dependente}
    f(x, y) =
    \begin{cases}
      c e^{-x-y}, & \text{se } 0 < x < y, \\[0.2cm]
      0, & \text{caso contrário},
    \end{cases}
  \end{equation}
  para alguma constante $c$. Encontre o valor de $c$ e determine se $X$ e $Y$
  são independentes.
  \begin{solution}
    A não negatividade da densidade implica $c \geqslant 0$. Como sua
    integral deve ser $1$, vamos integrar primeiro em relação a $y$,
    que varia de $x$ a $\infty$ na região $0<x<y$:
    \begin{align*}
      1
      &= \int_{-\infty}^{\infty}\int_{-\infty}^{\infty}
      f(x, y)\,dy\,dx
      \\[0.3cm]
      &= c\int_{0}^{\infty}\int_{x}^{\infty} e^{-x-y}\,dy\,dx
      \\[0.3cm]
      &= c\int_0^{\infty} e^{-2x}\,dx = \frac{c}{2},
    \end{align*}
    o que fornece $c=2$.

    Para verificar se $X$ e $Y$ são independentes, vamos comparar a
    probabilidade de $\{X>2,Y<1\}$ com o produto das probabilidades
    de $\{X>2\}$ e $\{Y<1\}$. Pela densidade conjunta,
    \[
      \mathbb{P}(X > 2, Y < 1) = \int_{2}^{\infty} \int_{-\infty}^{1}
      f(x, y)\,dy\,dx = 0,
    \]
    já que $f(x,y)=0$ quando $y<x$. Por outro lado,
    \begin{align*}
      \mathbb{P}(X > 2)
      &= \int_2^{\infty}\int_x^{\infty} 2e^{-x-y}\,dy\,dx
      = e^{-4} > 0,
      \\[0.3cm]
      \mathbb{P}(Y < 1)
      &= \int_0^1\int_0^y 2e^{-x-y}\,dx\,dy
      = (1-e^{-1})^2 > 0,
    \end{align*}
    de modo que
    \[
      \mathbb{P}(X > 2, Y < 1) \neq \mathbb{P}(X > 2)\mathbb{P}(Y < 1),
    \]
    mostrando que $X$ e $Y$ são dependentes.
  \end{solution}
\end{example}

\begin{exercise}\label{exercise:densidades-marginais-e-independencia}
  Sejam $X$ e $Y$ variáveis aleatórias com função densidade conjunta dada por
  \[
    f(x, y) =
    \begin{cases}
      cx, & \text{se } 0 < y < x < 1, \\[0.2cm]
      0, & \text{caso contrário}.
    \end{cases}
  \]
  Encontre o valor da constante $c$ e as funções densidade marginais de $X$
  e $Y$. As variáveis $X$ e $Y$ são independentes?
\end{exercise}

\begin{exercise}\label{exercise:independencia-tres-variaveis}
  As variáveis aleatórias $X$, $Y$ e $Z$ têm função densidade conjunta dada por
  \[
    f(x, y, z) =
    \begin{cases}
      8xyz, & \text{se } 0 < x,y,z < 1, \\[0.2cm]
      0, & \text{caso contrário}.
    \end{cases}
  \]
  As variáveis $X$, $Y$ e $Z$ são independentes? Calcule $\mathbb{P}(X > Y)$
  e $\mathbb{P}(Y > Z)$.
\end{exercise}

\section{Aplicações}\label{sec:aplicacoes-vetores-continuos}
Os conceitos de densidade conjunta, marginalização e independência permitem
investigar perguntas que aparecem em inteligência artificial. Vamos usá-los
para:
\begin{itemize}
\item construir sorteios com os pesos da função \emph{softmax};
\item identificar uma informação presente apenas na interação entre duas características;
\item examinar como a dimensão influencia a disponibilidade de amostras locais.
\end{itemize}
Antes do primeiro exemplo, estudaremos a distribuição exponencial, que
modelará os tempos de espera dos relógios usados no sorteio.

\subsection{A distribuição exponencial e a falta de memória}\label{subsec:exponencial-falta-memoria}

Uma variável aleatória não negativa pode representar o tempo de espera até
a ocorrência de um evento. No exemplo dos relógios da subseção seguinte,
cada variável representará o tempo até um relógio tocar. Uma hipótese
possível para esse tempo de espera é a seguinte: sabendo que o evento ainda
não ocorreu após $t$ unidades de tempo, a distribuição do tempo adicional
de espera é a mesma da espera inicial. Chamaremos essa propriedade de
\emph{falta de memória}.

Vamos definir a distribuição exponencial, verificar que ela possui essa
propriedade e mostrar que, entre as variáveis aleatórias estritamente
positivas quase certamente, ela é a única distribuição sem memória.
Depois relacionaremos essa propriedade à taxa de falha constante e a um
modelo de ocorrência de eventos em intervalos de tempo. A demonstração da
caracterização acompanha a apresentação de \cite[Seção~5.5]{Ross2010}.

\subsubsection{Definição e função de distribuição}

\begin{definition}[Distribuição exponencial]\label{def:distribuicao-exponencial}
Seja $\lambda>0$. Dizemos que uma variável aleatória $X$ tem
\emph{distribuição exponencial de parâmetro} $\lambda$, e escrevemos
$X\sim\mathrm{Exp}(\lambda)$, quando admite a função densidade
\[
  f_X(x)=
  \begin{cases}
    \lambda e^{-\lambda x}, & \text{se } x\geqslant 0,
    \\[0.2cm]
    0, & \text{se } x<0.
  \end{cases}
\]
O número $\lambda$ também é chamado de \emph{taxa} da distribuição.
\end{definition}

A função é não negativa, e sua integral sobre a reta é
\begin{align*}
  \int_{-\infty}^{\infty} f_X(x)\,dx
  &=\int_0^{\infty}\lambda e^{-\lambda x}\,dx
  =\left[-e^{-\lambda x}\right]_0^{\infty}=1.
\end{align*}
Portanto, ela é de fato uma densidade. Para $a<0$, temos $F_X(a)=0$;
para $a\geqslant0$, a função de distribuição é
\begin{equation}\label{eq:fda-exponencial}
  F_X(a)=\mathbb{P}(X\leqslant a)
  =\int_0^a\lambda e^{-\lambda x}\,dx=1-e^{-\lambda a}.
\end{equation}
Em particular, $\mathbb{P}(X>0)=1$. É conveniente registrar também a
\emph{função de cauda}, isto é, a probabilidade de $X$ exceder um valor
dado. Para $a\geqslant0$,
\begin{equation}\label{eq:cauda-exponencial}
  G(a)\equiv\mathbb{P}(X>a)=1-F_X(a)=e^{-\lambda a}.
\end{equation}
Essa função é não crescente e contínua, com $G(0)=1$.

\begin{figure}[H]
\centering
\begin{tikzpicture}
\begin{axis}[
  width=0.78\textwidth, height=6.2cm,
  axis lines=left,
  xlabel={$x$}, ylabel={$f_X$},
  ylabel style={rotate=-90},
  xmin=0, xmax=8, ymin=0, ymax=2.1,
  ytick={0,0.5,1,1.5,2},
  yticklabels={$0$,$\tfrac{1}{2}$,$1$,$\tfrac{3}{2}$,$2$},
  /pgf/number format/use comma,
  domain=0:8, samples=200,
  legend style={at={(0.97,0.97)},anchor=north east,draw=none,fill=none,
    row sep=4pt},
  every axis plot/.append style={line width=1.1pt},
]
\addplot[teal] {0.5*exp(-0.5*x)};
\addplot[muted] {1.0*exp(-1.0*x)};
\addplot[black] {2.0*exp(-2.0*x)};
\legend{$\lambda=\tfrac{1}{2}$,$\lambda=1$,$\lambda=2$}
\end{axis}
\end{tikzpicture}
\caption{Densidades exponenciais para três taxas. Quanto maior a taxa
$\lambda$, menor a probabilidade $\mathbb{P}(X>a)=e^{-\lambda a}$ para
cada $a>0$ fixado.}
\label{fig:densidades-exponenciais}
\end{figure}

\begin{example}[Média, variância e o significado da taxa]\label{ex:media-variancia-exponencial}
Calculemos $\mathbb{E}[X]$ para interpretar o parâmetro $\lambda$.
Integrando por partes com $u=x$ e $dv=\lambda e^{-\lambda x}\,dx$, de
modo que $v=-e^{-\lambda x}$, obtemos
\begin{align*}
  \mathbb{E}[X]
  &=\int_0^{\infty}x\lambda e^{-\lambda x}\,dx
  \\[0.3cm]
  &=\left[-xe^{-\lambda x}\right]_0^{\infty}
    +\int_0^{\infty}e^{-\lambda x}\,dx
  =\frac{1}{\lambda}.
\end{align*}
O termo de fronteira se anula porque $xe^{-\lambda x}\to0$ quando
$x\to\infty$. Aplicando integração por partes ao segundo momento,
com $u=x^2$ e o mesmo $dv$, temos
\[
  \mathbb{E}[X^2]
  =\left[-x^2e^{-\lambda x}\right]_0^{\infty}
    +2\int_0^{\infty}xe^{-\lambda x}\,dx
  =\frac{2}{\lambda}\mathbb{E}[X]=\frac{2}{\lambda^2}.
\]
Logo,
\[
  \operatorname{Var}(X)
  =\mathbb{E}[X^2]-\bigl(\mathbb{E}[X]\bigr)^2
  =\frac{1}{\lambda^2}.
\]
Assim, aumentar a taxa $\lambda$ reduz o tempo médio de espera
$1/\lambda$. Mais adiante veremos outra interpretação de $\lambda$:
ela determina a probabilidade de o evento ocorrer em um intervalo curto,
dado que ainda não ocorreu.
\end{example}

\subsubsection{A falta de memória}

Antes de enunciar a propriedade em tempo contínuo, recordemos seu análogo
discreto.

\begin{example}[A distribuição geométrica como análoga discreta]\label{ex:geometrica-sem-memoria}
Considere uma sequência de tentativas independentes, cada uma com
probabilidade de sucesso $p\in(0,1)$. Seja $N$ o número da tentativa em
que ocorre o primeiro sucesso. Então $N$ tem distribuição geométrica,
com a convenção
\[
  \mathbb{P}(N=n)=p(1-p)^{n-1},\qquad n\in\{1,2,\ldots\}.
\]
Para cada inteiro $k\geqslant0$,
\[
  \mathbb{P}(N>k)=(1-p)^k,
\]
pois $\{N>k\}$ é o evento de não haver sucesso nas $k$ primeiras
tentativas. Para inteiros $m,k\geqslant0$, a inclusão
$\{N>m+k\}\subseteq\{N>m\}$ dá
\[
  \mathbb{P}(N>m+k\mid N>m)
  =\frac{\mathbb{P}(N>m+k)}{\mathbb{P}(N>m)}
  =\frac{(1-p)^{m+k}}{(1-p)^m}
  =(1-p)^k=\mathbb{P}(N>k).
\]
Portanto, sob a condição de que as $m$ primeiras tentativas falharam,
$N-m$ tem a mesma distribuição que $N$. Essa é a propriedade de falta de
memória da geométrica. Sua cauda $(1-p)^k$ pode ser escrita como
$e^{-\lambda k}$, com $\lambda=-\log(1-p)>0$; a cauda exponencial
$e^{-\lambda t}$ estende essa forma aos tempos reais $t\geqslant0$.
\end{example}

\begin{definition}[Variável aleatória sem memória]\label{def:variavel-sem-memoria}
Uma variável aleatória $X$ satisfazendo a condição $\mathbb{P}(X\geqslant0)=1$ é
chamada de uma variável aleatória \emph{sem memória} se, para todos $s,t\geqslant0$ tais que
$\mathbb{P}(X>t)>0$,
\begin{equation}\label{eq:falta-memoria-condicional}
  \mathbb{P}(X>s+t\mid X>t)=\mathbb{P}(X>s).
\end{equation}
A condição $\mathbb{P}(X>t)>0$, na definição acima, é usada para garantir que a probabilidade condicional
esteja definida.
\end{definition}

No evento $\{X>t\}$, o tempo adicional de espera é $X-t$. Assim,
\eqref{eq:falta-memoria-condicional} também se escreve como
\[
  \mathbb{P}(X-t>s\mid X>t)=\mathbb{P}(X>s),\qquad s\geqslant0.
\]
Essa igualdade diz que a distribuição condicional do tempo restante,
após $t$ unidades de espera sem ocorrência do evento, é a mesma da espera
inicial. Para estudar a propriedade, vamos escrevê-la em termos da cauda
$G(x)=\mathbb{P}(X>x)$.

\begin{proposition}[Reformulação como equação funcional]\label{prop:falta-memoria-equacao}
Uma variável aleatória não negativa quase certamente $X$ é sem memória
se, e somente se, sua cauda $G(x)=\mathbb{P}(X>x)$ satisfaz
\begin{equation}\label{eq:equacao-funcional-cauda}
  G(s+t)=G(s)G(t),\qquad\forall s,t\geqslant0.
\end{equation}
\end{proposition}
\begin{proof}
Suponha primeiro que $X$ é sem memória e fixe $s,t\geqslant0$. Se
$G(t)>0$, a definição de probabilidade condicional e a inclusão
$\{X>s+t\}\subseteq\{X>t\}$ dão
\[
  G(s+t)=\mathbb{P}(X>s+t\mid X>t)\mathbb{P}(X>t)=G(s)G(t).
\]
Se $G(t)=0$, a monotonicidade da cauda dá
$0\leqslant G(s+t)\leqslant G(t)=0$. Portanto, a igualdade também vale
nesse caso. Reciprocamente, se \eqref{eq:equacao-funcional-cauda} vale
e $G(t)>0$, então
\[
  \mathbb{P}(X>s+t\mid X>t)=\frac{G(s+t)}{G(t)}=G(s)
  =\mathbb{P}(X>s).
\]
\end{proof}

\begin{proposition}[A exponencial é sem memória]\label{prop:exponencial-sem-memoria}
Se $X\sim\mathrm{Exp}(\lambda)$, então $X$ é sem memória.
\end{proposition}
\begin{proof}
Por \eqref{eq:cauda-exponencial}, sua cauda é $G(x)=e^{-\lambda x}$,
para $x\geqslant0$. Logo,
\[
  G(s+t)=e^{-\lambda(s+t)}=e^{-\lambda s}e^{-\lambda t}=G(s)G(t),
  \qquad s,t\geqslant0,
\]
e a afirmação segue da \Cref{prop:falta-memoria-equacao}.
\end{proof}

\subsubsection{A falta de memória caracteriza a exponencial}

A proposição anterior verifica que a exponencial é sem memória. Para
provar a recíproca, precisamos determinar as funções de cauda que resolvem
\eqref{eq:equacao-funcional-cauda}. A identidade
$e^{-\lambda(s+t)}=e^{-\lambda s}e^{-\lambda t}$ sugere a forma da
solução. Para demonstrar que essa forma é obrigatória, calcularemos
primeiro $G$ nos argumentos racionais, usando a equação funcional, e
depois estenderemos a argumentos em $\mathbb{R}$ usando a continuidade à direita.

\begin{theorem}[Caracterização pela falta de memória]\label{thm:caracterizacao-falta-memoria}
Seja $X$ uma variável aleatória com $\mathbb{P}(X>0)=1$. Se $X$ é sem
memória, então existe $\lambda>0$ tal que $X\sim\mathrm{Exp}(\lambda)$.
\end{theorem}
\begin{proof}
Escreva $G(x)=\mathbb{P}(X>x)$. Por hipótese, $G(0)=1$ e, pela
\Cref{prop:falta-memoria-equacao}, vale
\eqref{eq:equacao-funcional-cauda}.

\smallskip
\noindent\textbf{Passo 1 (determinar $G$ nos racionais).}
Seja $p=G(1)$. Como $G$ é não crescente e $G(0)=1$, temos
$0\leqslant p\leqslant1$. Vamos primeiro excluir os valores $0$ e $1$.

Se $p=1$, a iteração de \eqref{eq:equacao-funcional-cauda} dá
$G(n)=G(1)^n=1$ para todo inteiro $n\geqslant1$. Isso contradiz
$\lim_{n\to\infty}\mathbb{P}(X>n)=0$, pois $X$ assume valores reais.
Logo, $p<1$.

Se $p=0$, tomando $s=t=1/2$ em
\eqref{eq:equacao-funcional-cauda}, temos $G(1/2)^2=G(1)=0$.
Repetindo o argumento, $G(1/2^n)=0$ para todo inteiro $n\geqslant1$.
Os eventos $\{X>1/2^n\}$ crescem para $\{X>0\}$. Pela continuidade da
probabilidade em sequências crescentes,
\[
  G(0)=\mathbb{P}(X>0)
  =\lim_{n\to\infty}\mathbb{P}(X>1/2^n)=0,
\]
contrariando $G(0)=1$. Portanto, $0<p<1$.

Escolhemos então
\[
  \lambda=-\log p>0,\qquad\text{de modo que}\qquad p=e^{-\lambda}.
\]
Para inteiros $n\geqslant1$, a iteração da equação funcional dá
\[
  G(n)=G(1)^n=p^n=e^{-\lambda n}.
\]
Escrevendo $1$ como a soma de $n$ parcelas iguais a $1/n$ e aplicando a
mesma equação sucessivamente, obtemos
\[
  p=G(1)=G(1/n)^n,\qquad\text{logo}\qquad
  G(1/n)=p^{1/n}=e^{-\lambda/n}.
\]
Repetindo agora a parcela $1/n$ um total de $m$ vezes, para inteiros
$m,n\geqslant1$, segue que
\[
  G(m/n)=G(1/n)^m=p^{m/n}=e^{-\lambda m/n}.
\]
Juntamente com $G(0)=1$, isso determina $G(x)=e^{-\lambda x}$ para todo
racional $x\geqslant0$.

\smallskip
\noindent\textbf{Passo 2 (estender aos reais).}
Como $G=1-F_X$ e toda função de distribuição é contínua à direita, $G$
também é contínua à direita. Fixe $x\geqslant0$ e tome uma sequência de
racionais $x_k\downarrow x$. Então
\[
  G(x)=\lim_{k\to\infty}G(x_k)
  =\lim_{k\to\infty}e^{-\lambda x_k}=e^{-\lambda x}.
\]
Consequentemente, $F_X(x)=0$ para $x<0$ e
$F_X(x)=1-e^{-\lambda x}$ para $x\geqslant0$. Essa é a função de
distribuição calculada em \eqref{eq:fda-exponencial}. Em particular,
\[
  F_X(x)=\int_{-\infty}^{x}
  \lambda e^{-\lambda u}\mathds{1}_{[0,\infty)}(u)\,du,
  \qquad\forall x\in\mathbb{R},
\]
de modo que $X$ admite a densidade da
\Cref{def:distribuicao-exponencial} e $X\sim\mathrm{Exp}(\lambda)$.
\end{proof}

\begin{remark}[O caso degenerado]\label{rem:caso-degenerado}
Se $X=0$ quase certamente, então $\mathbb{P}(X>t)=0$ para todo
$t\geqslant0$. Nesse caso, a \Cref{def:variavel-sem-memoria} é satisfeita
de modo vacuamente verdadeiro: não há valor de $t$ para o qual seja
necessário verificar a igualdade condicional. A equação funcional também
vale, pois $G$ é identicamente nula em $[0,\infty)$.

Esse é o único caso sem memória que a hipótese
$\mathbb{P}(X>0)=1$ exclui entre as variáveis não negativas. De fato,
tomando $s=t=0$ em \eqref{eq:equacao-funcional-cauda}, temos
$G(0)=G(0)^2$, logo $G(0)\in\{0,1\}$. Se $G(0)=0$, então $X=0$ quase
certamente; se $G(0)=1$, aplica-se o
\Cref{thm:caracterizacao-falta-memoria}.
\end{remark}

\begin{remark}[O papel das propriedades da cauda]\label{rem:regularidade-equacao-funcional}
A equação funcional permite determinar $G$ nos racionais a partir de
$G(1)$. Para chegar aos reais, usamos a continuidade à direita, uma
propriedade que toda função de cauda já possui.

Alternativamente, podemos concluir essa etapa usando a monotonicidade de
$G$, sem recorrer à sua continuidade à direita. Para cada $x>0$,
escolhendo racionais não negativos $a_k\uparrow x$ e $b_k\downarrow x$,
temos
\[
  e^{-\lambda b_k}=G(b_k)\leqslant G(x)\leqslant G(a_k)=e^{-\lambda a_k}.
\]
Pela continuidade da função exponencial e pelo Teorema do Confronto,
segue que $G(x)=e^{-\lambda x}$. O caso $x=0$ já foi estabelecido.

Portanto, não foi necessário supor que $X$ tivesse densidade ou que $G$
fosse diferenciável. A existência da
densidade exponencial foi obtida como consequência da falta de memória.
\end{remark}

\subsubsection{Taxa de falha constante}

Seja $X$ não negativa quase certamente, com densidade $f_X$ contínua em
$[0,\infty)$ e cauda $G(t)=\mathbb{P}(X>t)>0$ para todo $t\geqslant0$.
Definimos sua \emph{taxa de falha}, ou \emph{taxa de risco}, por
\[
  h_X(t)=\frac{f_X(t)}{G(t)},\qquad t\geqslant0.
\]
Para entender essa razão, fixe um valor $t\geqslant0$. Para cada $\delta>0$, temos
\[
  \mathbb{P}(t<X\leqslant t+\delta\mid X>t)
  =\frac{1}{G(t)}\int_t^{t+\delta}f_X(u)\,du.
\]
Pela continuidade de $f_X$, dividindo por $\delta$ e fazendo
$\delta\downarrow0$, obtemos
\[
  \lim_{\delta\downarrow0}
  \frac{\mathbb{P}(t<X\leqslant t+\delta\mid X>t)}{\delta}
  =h_X(t).
\]
Assim, $h_X(t)$ mede a taxa de ocorrência do evento logo após $t$, sob a
condição de que ele ainda não ocorreu.

\begin{example}[Taxa de falha da exponencial]\label{ex:taxa-falha-exponencial}
Se $X\sim\mathrm{Exp}(\lambda)$, usando a densidade da
\Cref{def:distribuicao-exponencial}, temos
\[
  h_X(t)=\frac{\lambda e^{-\lambda t}}{e^{-\lambda t}}=\lambda,
  \qquad t\geqslant0.
\]
Portanto, essa taxa não depende do tempo já decorrido. Reciprocamente,
sob as hipóteses acima, suponha que $h_X(t)=\lambda>0$ para todo
$t\geqslant0$. Como $X$ tem densidade e é não negativa quase certamente,
$G(0)=1$. Além disso, para $t>0$,
\[
  G'(t)=-f_X(t)=-\lambda G(t).
\]
Para resolver essa equação, observe que
$\frac{d}{dt}(e^{\lambda t}G(t))=0$. Logo,
$e^{\lambda t}G(t)=G(0)=1$ e $G(t)=e^{-\lambda t}$, para
$t\geqslant0$. Assim, $X\sim\mathrm{Exp}(\lambda)$. Nessas condições,
a taxa de falha constante e a falta de memória caracterizam a mesma
distribuição.
\end{example}

\subsubsection{A exponencial como tempo de espera de eventos raros}

Podemos também obter a distribuição exponencial a partir de um modelo
para a contagem de eventos. A ideia é supor que as contagens em intervalos
disjuntos sejam independentes, que sua distribuição dependa apenas do
comprimento do intervalo e que, em um intervalo curto, a probabilidade
de uma ocorrência seja aproximadamente proporcional a esse comprimento.

\begin{proposition}[Tempo até o primeiro evento]\label{prop:tempo-primeiro-evento}
Sejam $N(t):\Omega\to\{0,1,2,\ldots\}$, para $t\geqslant0$, variáveis
aleatórias em $(\Omega,\mathcal{F},\mathbb{P})$, em que $N(t)$ representa
o número de eventos ocorridos no intervalo $(0,t]$. Suponha que, para
cada $\omega\in\Omega$, temos $N(0)(\omega)=0$ e a função
$t\mapsto N(t)(\omega)$ é não decrescente e contínua à direita.
Suponha ainda que:
\begin{enumerate}[label=\textnormal{(\roman*)}]
  \item as contagens em qualquer coleção finita de intervalos disjuntos
  sejam independentes;
  \item para quaisquer $a,h\geqslant0$, a contagem $N(a+h)-N(a)$ tenha
  a mesma distribuição que $N(h)$;
  \item exista $\lambda>0$ tal que, quando $h\downarrow0$,
  \[
    \mathbb{P}(N(h)=1)=\lambda h+o(h),
    \qquad \mathbb{P}(N(h)\geqslant2)=o(h).
  \]
\end{enumerate}
Aqui, $o(h)$ denota uma quantidade cuja razão por $h$ tende a zero
quando $h\downarrow0$. Então o tempo até o primeiro evento,
\[
  T=\inf\{t>0:N(t)\geqslant1\},
\]
com a convenção $\inf\varnothing=\infty$, é finito quase certamente e
tem distribuição $\mathrm{Exp}(\lambda)$.
\end{proposition}
\begin{proof}
O evento de o primeiro evento ocorrer após $t$ é o evento de não haver
ocorrências em $(0,t]$: $\{T>t\}=\{N(t)=0\}$.
Essa identidade decorre das propriedades das trajetórias de $N$.
Em particular, ela garante a mensurabilidade de $T$.

Para calcular essa probabilidade, fixe $t>0$ e divida $(0,t]$ em $n$
intervalos disjuntos de comprimento $t/n$. Pela hipótese (iii),
podemos escrever
\[
  \mathbb{P}(N(h)=0)=1-\lambda h+r(h),
  \qquad \lim_{h\downarrow0}\frac{r(h)}{h}=0.
\]
Pela hipótese (ii), essa é também a probabilidade de nenhuma ocorrência
em qualquer intervalo de comprimento $h$. Portanto, pela independência
das contagens em (i),
\[
  \mathbb{P}(T>t)
  =\mathbb{P}(N(t)=0)
  =\left(1-\frac{\lambda t}{n}+r(t/n)\right)^n.
\]
Como $nr(t/n)=t\,r(t/n)/(t/n)\to0$, o limite usual da função
exponencial dá
\[
  \lim_{n\to\infty}
  \left(1-\frac{\lambda t}{n}+r(t/n)\right)^n=e^{-\lambda t}.
\]
O lado esquerdo da igualdade anterior a esse limite é
$\mathbb{P}(T>t)$ para todo $n$. Logo,
$\mathbb{P}(T>t)=e^{-\lambda t}$ para todo $t>0$.
Fazendo $t\to\infty$ e $t\downarrow0$, respectivamente, concluímos
que $\mathbb{P}(T=\infty)=0$ e $\mathbb{P}(T>0)=1$.
Assim, $T$ tem a função de distribuição exponencial.
\end{proof}


\subsection{Corridas exponenciais e amostragem softmax}\label{subsec:aplicacoes-softmax}

\begin{example}[Relógios exponenciais, softmax e sorteios]\label{ex:corrida-exponenciais-softmax}
Considere dois relógios independentes $X\sim\mathrm{Exp}(\lambda)$ e
$Y\sim\mathrm{Exp}(\mu)$, com $\lambda,\mu>0$. O valor de cada uma destas
duas variáveis aleatórias representa o tempo que
o relógio correspondente leva para tocar.
\begin{figure}[H]
\centering
\begin{minipage}[t]{0.28\textwidth}
\centering
% TikZ conversion of figuras/despertador.svg (absolute M/L/C paths).
% Coordinates are normalized and flipped to the usual TikZ y-up axis.
% Resize by changing the x/y units; this size gives about 3 cm height.
\begin{tikzpicture}[x=0.0015cm,y=0.0015cm]
  \path[fill=black,nonzero rule]
    (213,752.113)
    .. controls (215.547,409.078) and (496.691,132.078) .. (839.711,134.625)
    .. controls (1182.762,137.156) and (1459.762,418.305) .. (1457.211,761.336)
    .. controls (1454.672,1104.371) and (1173.52,1381.371) .. (830.492,1378.832)
    .. controls (487.457,1376.281) and (210.457,1095.148) .. (213,752.113)
    -- cycle
    (1287.25,1160.09)
    -- (1231.621,1110.422)
    .. controls (1230.102,1109.059) and (1229.953,1106.711) .. (1231.293,1105.18)
    .. controls (1232.641,1103.66) and (1234.98,1103.5) .. (1236.531,1104.852)
    -- (1292.832,1153.73)
    .. controls (1328.672,1112.461) and (1359.043,1066.371) .. (1382.801,1016.48)
    -- (1315.332,984.441)
    .. controls (1313.48,983.551) and (1312.672,981.352) .. (1313.531,979.5)
    .. controls (1314.383,977.648) and (1316.582,976.84) .. (1318.441,977.699)
    -- (1386.363,1008.789)
    .. controls (1408.961,959.551) and (1425.281,906.832) .. (1433.992,851.582)
    -- (1360.062,839.949)
    .. controls (1358.031,839.629) and (1356.621,837.75) .. (1356.891,835.719)
    .. controls (1357.172,833.711) and (1359.02,832.309) .. (1361.051,832.602)
    -- (1435.121,843.18)
    .. controls (1438.973,816.367) and (1441.223,789.051) .. (1441.43,761.223)
    .. controls (1441.633,733.398) and (1439.793,706.035) .. (1436.34,679.176)
    -- (1362.113,688.66)
    .. controls (1360.09,688.918) and (1358.25,687.496) .. (1358,685.469)
    .. controls (1357.762,683.453) and (1359.211,681.578) .. (1361.23,681.285)
    -- (1435.332,670.758)
    .. controls (1427.43,615.379) and (1411.902,562.434) .. (1390.031,512.867)
    -- (1321.672,542.949)
    .. controls (1319.793,543.77) and (1317.613,542.922) .. (1316.781,541.066)
    .. controls (1315.961,539.211) and (1316.793,537.016) .. (1318.66,536.152)
    -- (1386.59,505.113)
    .. controls (1363.582,454.879) and (1333.891,408.34) .. (1298.672,366.559)
    -- (1241.652,414.594)
    .. controls (1240.082,415.918) and (1237.742,415.727) .. (1236.41,414.18)
    .. controls (1235.09,412.625) and (1235.281,410.289) .. (1236.82,408.945)
    -- (1293.18,360.117)
    .. controls (1257.281,318.691) and (1215.891,282.184) .. (1170.09,251.703)
    -- (1128.793,313.855)
    .. controls (1127.66,315.566) and (1125.363,316.031) .. (1123.66,314.922)
    .. controls (1121.953,313.812) and (1121.461,311.52) .. (1122.57,309.801)
    -- (1163,247.078)
    .. controls (1117.25,217.547) and (1067.34,193.961) .. (1014.262,177.512)
    -- (992.18,248.953)
    .. controls (991.582,250.91) and (989.512,252.031) .. (987.551,251.453)
    .. controls (985.602,250.883) and (984.492,248.828) .. (985.063,246.863)
    -- (1006.113,175.129)
    .. controls (954.531,159.934) and (900.141,151.344) .. (843.832,150.535)
    -- (842.762,225.391)
    .. controls (842.723,227.434) and (841.051,229.035) .. (839.012,228.973)
    .. controls (836.992,229.008) and (835.332,227.383) .. (835.332,225.328)
    -- (835.383,150.477)
    .. controls (779.07,150.438) and (724.543,158.223) .. (672.738,172.652)
    -- (692.73,244.695)
    .. controls (693.277,246.68) and (692.133,248.715) .. (690.168,249.25)
    .. controls (688.203,249.797) and (686.148,248.641) .. (685.578,246.68)
    -- (664.566,174.91)
    .. controls (611.242,190.574) and (560.984,213.418) .. (514.805,242.27)
    -- (554.293,305.582)
    .. controls (555.383,307.324) and (554.867,309.609) .. (553.137,310.699)
    .. controls (551.418,311.785) and (549.121,311.27) .. (548.012,309.543)
    -- (507.645,246.797)
    .. controls (461.41,276.59) and (419.473,312.484) .. (382.965,353.363)
    -- (438.598,403.023)
    .. controls (440.121,404.391) and (440.266,406.734) .. (438.922,408.258)
    .. controls (437.578,409.793) and (435.234,409.949) .. (433.688,408.605)
    -- (377.371,359.727)
    .. controls (341.535,400.984) and (311.172,447.074) .. (287.418,496.969)
    -- (354.883,529.012)
    .. controls (356.723,529.891) and (357.531,532.105) .. (356.688,533.953)
    .. controls (355.832,535.805) and (353.637,536.609) .. (351.766,535.754)
    -- (283.859,504.664)
    .. controls (261.258,553.906) and (244.941,606.617) .. (236.223,661.867)
    -- (310.152,673.496)
    .. controls (312.18,673.816) and (313.594,675.703) .. (313.328,677.73)
    .. controls (313.051,679.742) and (311.195,681.141) .. (309.16,680.855)
    -- (235.094,670.262)
    .. controls (231.25,697.086) and (228.996,724.402) .. (228.785,752.227)
    .. controls (228.582,780.051) and (230.426,807.414) .. (233.879,834.281)
    -- (308.105,824.793)
    .. controls (310.129,824.531) and (311.969,825.961) .. (312.207,827.98)
    .. controls (312.457,830) and (311.008,831.871) .. (308.977,832.172)
    -- (234.883,842.691)
    .. controls (242.773,898.07) and (258.312,951.02) .. (280.176,1000.59)
    -- (348.547,970.5)
    .. controls (350.426,969.68) and (352.609,970.52) .. (353.441,972.379)
    .. controls (354.258,974.242) and (353.426,976.43) .. (351.559,977.289)
    -- (283.625,1008.34)
    .. controls (306.637,1058.559) and (336.316,1105.109) .. (371.531,1146.891)
    -- (428.57,1098.859)
    .. controls (430.137,1097.531) and (432.48,1097.719) .. (433.793,1099.27)
    .. controls (435.117,1100.82) and (434.938,1103.16) .. (433.391,1104.512)
    -- (377.031,1153.332)
    .. controls (412.938,1194.762) and (454.332,1231.27) .. (500.125,1261.75)
    -- (541.414,1199.602)
    .. controls (542.555,1197.898) and (544.859,1197.422) .. (546.555,1198.531)
    .. controls (548.266,1199.641) and (548.742,1201.93) .. (547.633,1203.648)
    -- (507.215,1266.371)
    .. controls (552.961,1295.898) and (602.879,1319.492) .. (655.953,1335.941)
    -- (678.023,1264.5)
    .. controls (678.641,1262.539) and (680.699,1261.422) .. (682.664,1261.992)
    .. controls (684.609,1262.559) and (685.73,1264.609) .. (685.16,1266.59)
    -- (664.09,1338.32)
    .. controls (715.684,1353.52) and (770.078,1362.109) .. (826.383,1362.91)
    -- (827.453,1288.059)
    .. controls (827.48,1286.02) and (829.16,1284.422) .. (831.203,1284.48)
    .. controls (833.23,1284.441) and (834.883,1286.07) .. (834.891,1288.121)
    -- (834.84,1362.969)
    .. controls (891.152,1363) and (945.66,1355.23) .. (997.48,1340.789)
    -- (977.473,1268.742)
    .. controls (976.941,1266.77) and (978.082,1264.742) .. (980.051,1264.199)
    .. controls (982,1263.66) and (984.07,1264.809) .. (984.641,1266.77)
    -- (1005.652,1338.531)
    .. controls (1058.961,1322.879) and (1109.223,1300.02) .. (1155.402,1271.172)
    -- (1115.91,1207.871)
    .. controls (1114.832,1206.129) and (1115.352,1203.84) .. (1117.082,1202.75)
    .. controls (1118.801,1201.66) and (1121.082,1202.18) .. (1122.191,1203.898)
    -- (1162.57,1266.648)
    .. controls (1208.801,1236.859) and (1250.73,1200.969) .. (1287.25,1160.09);
  \path[fill=black,nonzero rule]
    (1642.66,1498.762)
    .. controls (1626.801,1566.82) and (1585.383,1624.629) .. (1526.031,1661.539)
    -- (1428.73,1722.051)
    .. controls (1386.602,1748.25) and (1338.23,1761.898) .. (1288.852,1761.531)
    .. controls (1198.473,1760.859) and (1116.031,1714.699) .. (1068.25,1637.84)
    -- (1027.562,1568.09)
    .. controls (1025.43,1564.41) and (1026.582,1559.672) .. (1030.211,1557.422)
    -- (1241.32,1426.129)
    -- (1203.43,1369.25)
    .. controls (1108.531,1426.531) and (999.363,1462.371) .. (882.711,1470.082)
    -- (882.352,1518.512)
    -- (967.23,1519.148)
    .. controls (971.602,1519.18) and (975.102,1522.73) .. (975.07,1527.09)
    -- (974.703,1576.391)
    .. controls (974.191,1644.422) and (918.441,1699.352) .. (850.422,1698.84)
    -- (805.82,1698.512)
    .. controls (737.805,1698.012) and (682.879,1642.262) .. (683.383,1574.242)
    -- (683.746,1524.93)
    .. controls (683.777,1520.57) and (687.332,1517.059) .. (691.695,1517.102)
    -- (776.562,1517.73)
    -- (776.93,1469.301)
    .. controls (659.891,1459.82) and (550.824,1422.059) .. (456.496,1362.93)
    -- (417.223,1420.012)
    -- (626.363,1554.422)
    .. controls (629.965,1556.742) and (631.062,1561.48) .. (628.855,1565.141)
    -- (587.043,1634.469)
    .. controls (538.223,1710.422) and (455.121,1755.359) .. (364.738,1754.68)
    .. controls (315.348,1754.309) and (267.195,1739.949) .. (225.457,1713.129)
    -- (129.07,1651.191)
    .. controls (70.27,1613.391) and (29.703,1554.98) .. (14.852,1486.691)
    .. controls (0,1418.398) and (12.633,1348.41) .. (50.297,1289.809)
    -- (91.996,1220.66)
    .. controls (93.098,1218.832) and (94.891,1217.531) .. (96.973,1217.051)
    .. controls (97.586,1216.898) and (98.207,1216.84) .. (98.82,1216.852)
    .. controls (100.297,1216.852) and (101.762,1217.281) .. (103.027,1218.09)
    -- (328.203,1362.809)
    -- (371.23,1300.262)
    .. controls (216.336,1167.871) and (118.563,970.66) .. (120.184,751.43)
    .. controls (121.477,577.914) and (184.801,419.133) .. (288.902,295.953)
    -- (232.887,54.953)
    .. controls (229.766,41.512) and (232.922,27.641) .. (241.535,16.875)
    .. controls (250.156,6.109) and (263.008,0) .. (276.793,0.102)
    -- (361.352,0.73)
    .. controls (365.012,0.75) and (368.16,3.285) .. (368.984,6.84)
    -- (409.629,182.547)
    .. controls (529.996,93.129) and (679.25,40.602) .. (840.41,41.805)
    .. controls (1001.57,42.996) and (1150.02,97.723) .. (1269.051,188.918)
    -- (1312.293,13.836)
    .. controls (1313.16,10.289) and (1316.363,7.812) .. (1320.012,7.832)
    -- (1404.582,8.457)
    .. controls (1418.371,8.562) and (1431.133,14.863) .. (1439.582,25.762)
    .. controls (1448.043,36.648) and (1450.98,50.582) .. (1447.672,63.949)
    -- (1388.082,304.098)
    .. controls (1490.34,428.809) and (1551.32,588.516) .. (1550.031,762.031)
    .. controls (1548.402,981.801) and (1447.242,1177.949) .. (1289.75,1308.031)
    -- (1331.18,1370.242)
    -- (1558.492,1228.891)
    .. controls (1559.762,1228.09) and (1561.23,1227.691) .. (1562.711,1227.699)
    .. controls (1563.332,1227.699) and (1563.953,1227.781) .. (1564.551,1227.93)
    .. controls (1566.621,1228.449) and (1568.402,1229.781) .. (1569.473,1231.621)
    -- (1610.043,1301.18)
    .. controls (1646.953,1360.52) and (1658.543,1430.691) .. (1642.66,1498.762)
    -- cycle
    (699.168,1574.352)
    .. controls (698.723,1633.68) and (746.625,1682.289) .. (805.953,1682.73)
    -- (850.543,1683.059)
    .. controls (909.863,1683.492) and (958.473,1635.602) .. (958.91,1576.281)
    -- (959.223,1534.859)
    -- (874.352,1534.23)
    -- (784.34,1533.559)
    -- (699.461,1532.949)
    -- cycle
    (792.352,1517.852)
    -- (866.57,1518.391)
    -- (866.93,1470.969)
    .. controls (854.621,1471.5) and (842.25,1471.73) .. (829.813,1471.648)
    .. controls (817.352,1471.551) and (804.992,1471.129) .. (792.703,1470.41)
    -- cycle
    (101.324,1235.77)
    -- (63.688,1298.148)
    .. controls (28.184,1353.391) and (16.324,1419.16) .. (30.285,1483.332)
    .. controls (44.246,1547.5) and (82.352,1602.398) .. (137.594,1637.898)
    -- (233.992,1699.852)
    .. controls (273.211,1725.051) and (318.461,1738.559) .. (364.848,1738.898)
    .. controls (449.797,1739.531) and (527.891,1697.301) .. (573.641,1626.129)
    -- (611.387,1563.559)
    -- cycle
    (341.48,1371.34)
    -- (403.941,1411.48)
    -- (443.203,1354.41)
    .. controls (422.488,1340.781) and (402.531,1326.082) .. (383.387,1310.43)
    -- cycle
    (355.012,16.465)
    -- (276.68,15.887)
    .. controls (267.75,15.816) and (259.441,19.766) .. (253.859,26.742)
    .. controls (248.281,33.703) and (246.25,42.688) .. (248.273,51.375)
    -- (301.688,281.246)
    .. controls (330.371,249.105) and (361.883,219.574) .. (395.855,193.039)
    -- cycle
    (1432.34,60.16)
    .. controls (1434.5,51.492) and (1432.59,42.484) .. (1427.113,35.441)
    .. controls (1421.641,28.395) and (1413.391,24.312) .. (1404.461,24.242)
    -- (1326.121,23.664)
    -- (1282.672,199.617)
    .. controls (1316.25,226.648) and (1347.32,256.648) .. (1375.52,289.203)
    -- cycle
    (1518.461,761.797)
    .. controls (1521.262,384.984) and (1216.98,76.16) .. (840.172,73.371)
    .. controls (463.363,70.582) and (154.551,374.859) .. (151.758,751.656)
    .. controls (148.957,1128.461) and (453.238,1437.289) .. (830.031,1440.082)
    .. controls (1206.84,1442.871) and (1515.672,1138.602) .. (1518.461,761.797)
    -- cycle
    (1277.41,1317.992)
    .. controls (1258,1333.309) and (1237.812,1347.68) .. (1216.871,1360.949)
    -- (1254.742,1417.781)
    -- (1317.781,1378.582)
    -- cycle
    (1596.52,1309.309)
    -- (1559.93,1246.582)
    -- (1333.043,1387.691)
    -- (1287.742,1415.852)
    -- (1045.07,1566.781)
    -- (1081.77,1629.699)
    .. controls (1126.562,1701.742) and (1204.031,1745.121) .. (1288.973,1745.75)
    .. controls (1335.363,1746.102) and (1380.801,1733.262) .. (1420.391,1708.648)
    -- (1517.703,1648.141)
    .. controls (1573.473,1613.461) and (1612.383,1559.129) .. (1627.293,1495.172)
    .. controls (1642.203,1431.211) and (1631.312,1365.281) .. (1596.52,1309.309);
  % Ponteiros no mesmo estilo de traço e ponta usado no despertador vermelho.
  \draw[black,line width=1.1pt,line cap=round,
    -{Stealth[length=5pt,width=4pt]}] (835,757) -- (1363.223,654.453);
  \draw[black,line width=1.4pt,line cap=round,
    -{Stealth[length=5pt,width=4pt]}] (835,757) -- (909.402,1203.141);
  \fill[black] (835,757) circle[radius=13];
\end{tikzpicture}

\\[0.3cm]
{\small $X\sim \mathrm{Exp}(\lambda)$}
\end{minipage}\hspace{2.6cm}%
\begin{minipage}[t]{0.28\textwidth}
\centering
% TikZ conversion of figuras/despertador.svg (absolute M/L/C paths).
% Coordinates are normalized and flipped to the usual TikZ y-up axis.
% Resize by changing the x/y units; this size gives about 3 cm height.
\begin{tikzpicture}[x=0.0015cm,y=0.0015cm]
  \path[fill=red!80!black,nonzero rule]
    (213,752.113)
    .. controls (215.547,409.078) and (496.691,132.078) .. (839.711,134.625)
    .. controls (1182.762,137.156) and (1459.762,418.305) .. (1457.211,761.336)
    .. controls (1454.672,1104.371) and (1173.52,1381.371) .. (830.492,1378.832)
    .. controls (487.457,1376.281) and (210.457,1095.148) .. (213,752.113)
    -- cycle
    (1287.25,1160.09)
    -- (1231.621,1110.422)
    .. controls (1230.102,1109.059) and (1229.953,1106.711) .. (1231.293,1105.18)
    .. controls (1232.641,1103.66) and (1234.98,1103.5) .. (1236.531,1104.852)
    -- (1292.832,1153.73)
    .. controls (1328.672,1112.461) and (1359.043,1066.371) .. (1382.801,1016.48)
    -- (1315.332,984.441)
    .. controls (1313.48,983.551) and (1312.672,981.352) .. (1313.531,979.5)
    .. controls (1314.383,977.648) and (1316.582,976.84) .. (1318.441,977.699)
    -- (1386.363,1008.789)
    .. controls (1408.961,959.551) and (1425.281,906.832) .. (1433.992,851.582)
    -- (1360.062,839.949)
    .. controls (1358.031,839.629) and (1356.621,837.75) .. (1356.891,835.719)
    .. controls (1357.172,833.711) and (1359.02,832.309) .. (1361.051,832.602)
    -- (1435.121,843.18)
    .. controls (1438.973,816.367) and (1441.223,789.051) .. (1441.43,761.223)
    .. controls (1441.633,733.398) and (1439.793,706.035) .. (1436.34,679.176)
    -- (1362.113,688.66)
    .. controls (1360.09,688.918) and (1358.25,687.496) .. (1358,685.469)
    .. controls (1357.762,683.453) and (1359.211,681.578) .. (1361.23,681.285)
    -- (1435.332,670.758)
    .. controls (1427.43,615.379) and (1411.902,562.434) .. (1390.031,512.867)
    -- (1321.672,542.949)
    .. controls (1319.793,543.77) and (1317.613,542.922) .. (1316.781,541.066)
    .. controls (1315.961,539.211) and (1316.793,537.016) .. (1318.66,536.152)
    -- (1386.59,505.113)
    .. controls (1363.582,454.879) and (1333.891,408.34) .. (1298.672,366.559)
    -- (1241.652,414.594)
    .. controls (1240.082,415.918) and (1237.742,415.727) .. (1236.41,414.18)
    .. controls (1235.09,412.625) and (1235.281,410.289) .. (1236.82,408.945)
    -- (1293.18,360.117)
    .. controls (1257.281,318.691) and (1215.891,282.184) .. (1170.09,251.703)
    -- (1128.793,313.855)
    .. controls (1127.66,315.566) and (1125.363,316.031) .. (1123.66,314.922)
    .. controls (1121.953,313.812) and (1121.461,311.52) .. (1122.57,309.801)
    -- (1163,247.078)
    .. controls (1117.25,217.547) and (1067.34,193.961) .. (1014.262,177.512)
    -- (992.18,248.953)
    .. controls (991.582,250.91) and (989.512,252.031) .. (987.551,251.453)
    .. controls (985.602,250.883) and (984.492,248.828) .. (985.063,246.863)
    -- (1006.113,175.129)
    .. controls (954.531,159.934) and (900.141,151.344) .. (843.832,150.535)
    -- (842.762,225.391)
    .. controls (842.723,227.434) and (841.051,229.035) .. (839.012,228.973)
    .. controls (836.992,229.008) and (835.332,227.383) .. (835.332,225.328)
    -- (835.383,150.477)
    .. controls (779.07,150.438) and (724.543,158.223) .. (672.738,172.652)
    -- (692.73,244.695)
    .. controls (693.277,246.68) and (692.133,248.715) .. (690.168,249.25)
    .. controls (688.203,249.797) and (686.148,248.641) .. (685.578,246.68)
    -- (664.566,174.91)
    .. controls (611.242,190.574) and (560.984,213.418) .. (514.805,242.27)
    -- (554.293,305.582)
    .. controls (555.383,307.324) and (554.867,309.609) .. (553.137,310.699)
    .. controls (551.418,311.785) and (549.121,311.27) .. (548.012,309.543)
    -- (507.645,246.797)
    .. controls (461.41,276.59) and (419.473,312.484) .. (382.965,353.363)
    -- (438.598,403.023)
    .. controls (440.121,404.391) and (440.266,406.734) .. (438.922,408.258)
    .. controls (437.578,409.793) and (435.234,409.949) .. (433.688,408.605)
    -- (377.371,359.727)
    .. controls (341.535,400.984) and (311.172,447.074) .. (287.418,496.969)
    -- (354.883,529.012)
    .. controls (356.723,529.891) and (357.531,532.105) .. (356.688,533.953)
    .. controls (355.832,535.805) and (353.637,536.609) .. (351.766,535.754)
    -- (283.859,504.664)
    .. controls (261.258,553.906) and (244.941,606.617) .. (236.223,661.867)
    -- (310.152,673.496)
    .. controls (312.18,673.816) and (313.594,675.703) .. (313.328,677.73)
    .. controls (313.051,679.742) and (311.195,681.141) .. (309.16,680.855)
    -- (235.094,670.262)
    .. controls (231.25,697.086) and (228.996,724.402) .. (228.785,752.227)
    .. controls (228.582,780.051) and (230.426,807.414) .. (233.879,834.281)
    -- (308.105,824.793)
    .. controls (310.129,824.531) and (311.969,825.961) .. (312.207,827.98)
    .. controls (312.457,830) and (311.008,831.871) .. (308.977,832.172)
    -- (234.883,842.691)
    .. controls (242.773,898.07) and (258.312,951.02) .. (280.176,1000.59)
    -- (348.547,970.5)
    .. controls (350.426,969.68) and (352.609,970.52) .. (353.441,972.379)
    .. controls (354.258,974.242) and (353.426,976.43) .. (351.559,977.289)
    -- (283.625,1008.34)
    .. controls (306.637,1058.559) and (336.316,1105.109) .. (371.531,1146.891)
    -- (428.57,1098.859)
    .. controls (430.137,1097.531) and (432.48,1097.719) .. (433.793,1099.27)
    .. controls (435.117,1100.82) and (434.938,1103.16) .. (433.391,1104.512)
    -- (377.031,1153.332)
    .. controls (412.938,1194.762) and (454.332,1231.27) .. (500.125,1261.75)
    -- (541.414,1199.602)
    .. controls (542.555,1197.898) and (544.859,1197.422) .. (546.555,1198.531)
    .. controls (548.266,1199.641) and (548.742,1201.93) .. (547.633,1203.648)
    -- (507.215,1266.371)
    .. controls (552.961,1295.898) and (602.879,1319.492) .. (655.953,1335.941)
    -- (678.023,1264.5)
    .. controls (678.641,1262.539) and (680.699,1261.422) .. (682.664,1261.992)
    .. controls (684.609,1262.559) and (685.73,1264.609) .. (685.16,1266.59)
    -- (664.09,1338.32)
    .. controls (715.684,1353.52) and (770.078,1362.109) .. (826.383,1362.91)
    -- (827.453,1288.059)
    .. controls (827.48,1286.02) and (829.16,1284.422) .. (831.203,1284.48)
    .. controls (833.23,1284.441) and (834.883,1286.07) .. (834.891,1288.121)
    -- (834.84,1362.969)
    .. controls (891.152,1363) and (945.66,1355.23) .. (997.48,1340.789)
    -- (977.473,1268.742)
    .. controls (976.941,1266.77) and (978.082,1264.742) .. (980.051,1264.199)
    .. controls (982,1263.66) and (984.07,1264.809) .. (984.641,1266.77)
    -- (1005.652,1338.531)
    .. controls (1058.961,1322.879) and (1109.223,1300.02) .. (1155.402,1271.172)
    -- (1115.91,1207.871)
    .. controls (1114.832,1206.129) and (1115.352,1203.84) .. (1117.082,1202.75)
    .. controls (1118.801,1201.66) and (1121.082,1202.18) .. (1122.191,1203.898)
    -- (1162.57,1266.648)
    .. controls (1208.801,1236.859) and (1250.73,1200.969) .. (1287.25,1160.09);
  \path[fill=red!80!black,nonzero rule]
    (1642.66,1498.762)
    .. controls (1626.801,1566.82) and (1585.383,1624.629) .. (1526.031,1661.539)
    -- (1428.73,1722.051)
    .. controls (1386.602,1748.25) and (1338.23,1761.898) .. (1288.852,1761.531)
    .. controls (1198.473,1760.859) and (1116.031,1714.699) .. (1068.25,1637.84)
    -- (1027.562,1568.09)
    .. controls (1025.43,1564.41) and (1026.582,1559.672) .. (1030.211,1557.422)
    -- (1241.32,1426.129)
    -- (1203.43,1369.25)
    .. controls (1108.531,1426.531) and (999.363,1462.371) .. (882.711,1470.082)
    -- (882.352,1518.512)
    -- (967.23,1519.148)
    .. controls (971.602,1519.18) and (975.102,1522.73) .. (975.07,1527.09)
    -- (974.703,1576.391)
    .. controls (974.191,1644.422) and (918.441,1699.352) .. (850.422,1698.84)
    -- (805.82,1698.512)
    .. controls (737.805,1698.012) and (682.879,1642.262) .. (683.383,1574.242)
    -- (683.746,1524.93)
    .. controls (683.777,1520.57) and (687.332,1517.059) .. (691.695,1517.102)
    -- (776.562,1517.73)
    -- (776.93,1469.301)
    .. controls (659.891,1459.82) and (550.824,1422.059) .. (456.496,1362.93)
    -- (417.223,1420.012)
    -- (626.363,1554.422)
    .. controls (629.965,1556.742) and (631.062,1561.48) .. (628.855,1565.141)
    -- (587.043,1634.469)
    .. controls (538.223,1710.422) and (455.121,1755.359) .. (364.738,1754.68)
    .. controls (315.348,1754.309) and (267.195,1739.949) .. (225.457,1713.129)
    -- (129.07,1651.191)
    .. controls (70.27,1613.391) and (29.703,1554.98) .. (14.852,1486.691)
    .. controls (0,1418.398) and (12.633,1348.41) .. (50.297,1289.809)
    -- (91.996,1220.66)
    .. controls (93.098,1218.832) and (94.891,1217.531) .. (96.973,1217.051)
    .. controls (97.586,1216.898) and (98.207,1216.84) .. (98.82,1216.852)
    .. controls (100.297,1216.852) and (101.762,1217.281) .. (103.027,1218.09)
    -- (328.203,1362.809)
    -- (371.23,1300.262)
    .. controls (216.336,1167.871) and (118.563,970.66) .. (120.184,751.43)
    .. controls (121.477,577.914) and (184.801,419.133) .. (288.902,295.953)
    -- (232.887,54.953)
    .. controls (229.766,41.512) and (232.922,27.641) .. (241.535,16.875)
    .. controls (250.156,6.109) and (263.008,0) .. (276.793,0.102)
    -- (361.352,0.73)
    .. controls (365.012,0.75) and (368.16,3.285) .. (368.984,6.84)
    -- (409.629,182.547)
    .. controls (529.996,93.129) and (679.25,40.602) .. (840.41,41.805)
    .. controls (1001.57,42.996) and (1150.02,97.723) .. (1269.051,188.918)
    -- (1312.293,13.836)
    .. controls (1313.16,10.289) and (1316.363,7.812) .. (1320.012,7.832)
    -- (1404.582,8.457)
    .. controls (1418.371,8.562) and (1431.133,14.863) .. (1439.582,25.762)
    .. controls (1448.043,36.648) and (1450.98,50.582) .. (1447.672,63.949)
    -- (1388.082,304.098)
    .. controls (1490.34,428.809) and (1551.32,588.516) .. (1550.031,762.031)
    .. controls (1548.402,981.801) and (1447.242,1177.949) .. (1289.75,1308.031)
    -- (1331.18,1370.242)
    -- (1558.492,1228.891)
    .. controls (1559.762,1228.09) and (1561.23,1227.691) .. (1562.711,1227.699)
    .. controls (1563.332,1227.699) and (1563.953,1227.781) .. (1564.551,1227.93)
    .. controls (1566.621,1228.449) and (1568.402,1229.781) .. (1569.473,1231.621)
    -- (1610.043,1301.18)
    .. controls (1646.953,1360.52) and (1658.543,1430.691) .. (1642.66,1498.762)
    -- cycle
    (699.168,1574.352)
    .. controls (698.723,1633.68) and (746.625,1682.289) .. (805.953,1682.73)
    -- (850.543,1683.059)
    .. controls (909.863,1683.492) and (958.473,1635.602) .. (958.91,1576.281)
    -- (959.223,1534.859)
    -- (874.352,1534.23)
    -- (784.34,1533.559)
    -- (699.461,1532.949)
    -- cycle
    (792.352,1517.852)
    -- (866.57,1518.391)
    -- (866.93,1470.969)
    .. controls (854.621,1471.5) and (842.25,1471.73) .. (829.813,1471.648)
    .. controls (817.352,1471.551) and (804.992,1471.129) .. (792.703,1470.41)
    -- cycle
    (101.324,1235.77)
    -- (63.688,1298.148)
    .. controls (28.184,1353.391) and (16.324,1419.16) .. (30.285,1483.332)
    .. controls (44.246,1547.5) and (82.352,1602.398) .. (137.594,1637.898)
    -- (233.992,1699.852)
    .. controls (273.211,1725.051) and (318.461,1738.559) .. (364.848,1738.898)
    .. controls (449.797,1739.531) and (527.891,1697.301) .. (573.641,1626.129)
    -- (611.387,1563.559)
    -- cycle
    (341.48,1371.34)
    -- (403.941,1411.48)
    -- (443.203,1354.41)
    .. controls (422.488,1340.781) and (402.531,1326.082) .. (383.387,1310.43)
    -- cycle
    (355.012,16.465)
    -- (276.68,15.887)
    .. controls (267.75,15.816) and (259.441,19.766) .. (253.859,26.742)
    .. controls (248.281,33.703) and (246.25,42.688) .. (248.273,51.375)
    -- (301.688,281.246)
    .. controls (330.371,249.105) and (361.883,219.574) .. (395.855,193.039)
    -- cycle
    (1432.34,60.16)
    .. controls (1434.5,51.492) and (1432.59,42.484) .. (1427.113,35.441)
    .. controls (1421.641,28.395) and (1413.391,24.312) .. (1404.461,24.242)
    -- (1326.121,23.664)
    -- (1282.672,199.617)
    .. controls (1316.25,226.648) and (1347.32,256.648) .. (1375.52,289.203)
    -- cycle
    (1518.461,761.797)
    .. controls (1521.262,384.984) and (1216.98,76.16) .. (840.172,73.371)
    .. controls (463.363,70.582) and (154.551,374.859) .. (151.758,751.656)
    .. controls (148.957,1128.461) and (453.238,1437.289) .. (830.031,1440.082)
    .. controls (1206.84,1442.871) and (1515.672,1138.602) .. (1518.461,761.797)
    -- cycle
    (1277.41,1317.992)
    .. controls (1258,1333.309) and (1237.812,1347.68) .. (1216.871,1360.949)
    -- (1254.742,1417.781)
    -- (1317.781,1378.582)
    -- cycle
    (1596.52,1309.309)
    -- (1559.93,1246.582)
    -- (1333.043,1387.691)
    -- (1287.742,1415.852)
    -- (1045.07,1566.781)
    -- (1081.77,1629.699)
    .. controls (1126.562,1701.742) and (1204.031,1745.121) .. (1288.973,1745.75)
    .. controls (1335.363,1746.102) and (1380.801,1733.262) .. (1420.391,1708.648)
    -- (1517.703,1648.141)
    .. controls (1573.473,1613.461) and (1612.383,1559.129) .. (1627.293,1495.172)
    .. controls (1642.203,1431.211) and (1631.312,1365.281) .. (1596.52,1309.309);
  % Ponteiros marcando 10h10.
  \draw[red!80!black,line width=1.1pt,line cap=round,
    -{Stealth[length=5pt,width=4pt]}] (835,757) -- (1233,987);
  \draw[red!80!black,line width=1.4pt,line cap=round,
    -{Stealth[length=5pt,width=4pt]}] (835,757) -- (589,929);
  \fill[red!80!black] (835,757) circle[radius=13];
\end{tikzpicture}

\\[0.3cm]
{\small $Y\sim \mathrm{Exp}(\mu)$}
\end{minipage}
\label{fig:despertadores-tikz}
\end{figure}


Pela independência e pela forma das densidades exponenciais, uma densidade
conjunta é
\[
  f_{X,Y}(x,y)
  = \lambda e^{-\lambda x}\,\mu e^{-\mu y}
  \mathds{1}_{(0,\infty)^2}(x,y).
\]
De fato, a integral desse produto sobre
$(-\infty,x]\times(-\infty,y]$ é $F_X(x)F_Y(y)=F_{X,Y}(x,y)$, como exige
a definição de densidade conjunta.
Vamos primeiro calcular a probabilidade de que o relógio $X$ toque antes de
$Y$. No plano, o evento $\{X<Y\}$ é a região acima da diagonal $y=x$ no
quadrante positivo, como ilustra a Figura~\ref{fig:regiao-x-menor-y}.
\begin{figure}[H]
\centering
\begin{tikzpicture}[scale=0.9,>=Stealth]
  \fill[teal!15] (0,0) -- (0,4) -- (4,4) -- cycle;
  \draw[->,thick] (-1,0) -- (4.5,0) node[below] {$x$};
  \draw[->,thick] (0,-0.7) -- (0,4.5) node[left] {$y$};
  \draw[thick,dashed,muted] (0,0) -- (4,4)
    node[pos=0.72,below right=2pt,text=text] {$y=x$};
  % Seção vertical em x=x_0, com as projeções sobre os dois eixos.
  \draw[red!75,dashed] (0,2.2) -- (2.2,2.2);
  \draw[red!75,dashed] (2.2,0) -- (2.2,2.2);
  \draw[red!75,thick] (2.2,2.25) -- (2.2,4);
  \fill[red!75] (2.2,2.2) circle (1.5pt);
  \node at (-0.35,2.2) {$x_0$};
  \node at (2.2,-0.35) {$x_0$};
  \node[align=center] at (1.15,3.15) {$x<y$};
  \node[below left=2pt] at (0,0) {$0$};
\end{tikzpicture}
\captionsetup{width=\textwidth}
\caption{A região sombreada representa $\mathcal{S}\equiv \{(x,y)\in(0,\infty)^2:x<y\}$.}
\label{fig:regiao-x-menor-y}
\end{figure}
Para descrever essa região por seções verticais, fixemos $x=x_0>0$. A seção
indicada em vermelho na figura mostra que, para esse valor de $x$, a coordenada
$y$ varia de $x_0$ até $\infty$. Portanto
\begin{align*}
  \mathbb{P}(X<Y)
  =
  \mathbb{P}\bigl( (X,Y)\in\mathcal{S} \bigr)
  =
  \iint_{\mathcal{S}} f_{X,Y}(x,y)\, dx\,dy
  &= \int_0^\infty\int_x^\infty
  \lambda\mu e^{-\lambda x-\mu y}\,dy\,dx
  \\[0.3cm]
  &= \int_0^\infty \lambda e^{-\lambda x}e^{-\mu x}\,dx
  = \frac{\lambda}{\lambda+\mu}.
\end{align*}
Trocando os papéis dos relógios, obtemos
$\mathbb{P}(Y<X)=\mu/(\lambda+\mu)$. O empate corresponde à diagonal
$\{(x,y):x=y\}$, que tem área nula no plano. Como $(X,Y)$ tem densidade,
segue que $\mathbb{P}(X=Y)=0$. Assim, quase certamente há um único relógio
que toca primeiro.

O cálculo sugere que a taxa de um relógio determina sua chance de vencer.
Agora queremos acompanhar também o instante do primeiro toque. Consideremos
$n\geqslant2$ relógios independentes $X_1,\ldots,X_n$, com
$X_i\sim\mathrm{Exp}(\lambda_i)$ e $\lambda_i>0$. Para registrar o instante
e o relógio que toca primeiro, definimos
\[
  T=\min_{1\leqslant j\leqslant n}X_j,
  \qquad
  J=\min\{i\in\{1,\ldots,n\}:X_i=T\},
  \qquad
  \Lambda=\sum_{j=1}^n\lambda_j.
\]
Assim, $T$ é o instante do primeiro toque, e $J$ indica qual relógio o
produziu. A definição também cobre o caso de empate: escolhemos o menor índice
entre os relógios que tocam juntos. Vamos ver que essa convenção só afeta um
evento de probabilidade zero. Fixemos um par $i<j$. Pela independência, a
densidade conjunta de $(X_i,X_j)$ é
$f_{X_i,X_j}(x,y)=\lambda_i\lambda_j e^{-\lambda_i x-\lambda_j y}$, para
$x,y>0$. O empate entre esses relógios corresponde à diagonal $x=y$. Como a
probabilidade é a integral da densidade nessa região, e a diagonal tem área
nula, temos $\mathbb{P}(X_i=X_j)=0$. Se houver empate no primeiro toque,
algum par de relógios terá tempos iguais. Há apenas finitos pares, e a
probabilidade da união dos eventos $\{X_i=X_j\}$ é no máximo a soma das
probabilidades individuais; portanto, também é zero.

Para calcular a lei conjunta de $T$ e $J$, fixemos $i$ e consideremos a
probabilidade de o relógio $i$ ser o primeiro a tocar, com esse toque ocorrendo
depois de $t\geqslant0$.
Como os empates têm probabilidade zero, podemos integrar sobre a região em que
$i$ toca estritamente antes dos demais. Se esse toque ocorre no instante
$u>t$, todas as outras coordenadas $x_j$ devem satisfazer $x_j>u$. Integramos
primeiro nessas coordenadas e depois em $u$:
\begin{align}\label{eq-aux1-ex-corrida-despertadores}
  \mathbb{P}(T>t, J=i)
  &= \int_t^\infty
  \lambda_i e^{-\lambda_i u}
  \prod_{j\ne i}\left(\int_u^\infty
  \lambda_j e^{-\lambda_j x_j}\,dx_j\right)du
  \nonumber\\[0.3cm]
  &= \int_t^\infty \lambda_i e^{-\Lambda u}\,du
  = \frac{\lambda_i}{\Lambda}e^{-\Lambda t}.
\end{align}
Cada integral interna integra a densidade do relógio $j$ nos valores $x_j>u$,
enquanto a integral externa percorre os instantes $u>t$ em que o relógio $i$
pode tocar primeiro.

Somando sobre $i\in\{1,\ldots,n\}$, usamos a aditividade de $\mathbb{P}$ para obter
\begin{align}\label{eq-aux2-ex-corrida-despertadores} 
  \mathbb{P}(T>t)=e^{-\Lambda t},
\end{align}
portanto $T\sim\mathrm{Exp}(\Lambda)$. Fazendo $t=0$ em \eqref{eq-aux1-ex-corrida-despertadores},
obtemos $\mathbb{P}(T>0, J=i)=\lambda_i/\Lambda$. Como $T>0$ quase
certamente, temos $\mathbb{P}(T>0,J=i)=\mathbb{P}(J=i)$, e segue que
\begin{align}\label{eq-aux3-ex-corrida-despertadores}
\mathbb{P}(J=i)=\frac{\lambda_i}{\Lambda}.
\end{align}
Além disso, segue das identidades 
\eqref{eq-aux1-ex-corrida-despertadores}, 
\eqref{eq-aux2-ex-corrida-despertadores} 
e 
\eqref{eq-aux3-ex-corrida-despertadores}
que, para todo $t\geqslant0$,
\begin{align}\label{eq-aux4-ex-corrida-despertadores}
  \mathbb{P}(T>t,J=i)
  =\mathbb{P}(T>t)\,\mathbb{P}(J=i).
\end{align}

Fixado $k\in\mathbb{R}$, somando a igualdade acima sobre os índices $i\in\{1,\ldots,n\}$ satisfazendo
$i\leqslant k$, obtemos
\[
  \mathbb{P}(T>t,J\leqslant k)
  = \sum_{1\leqslant i\leqslant k} \mathbb{P}(T>t,J=i)
  = \mathbb{P}(T>t)\sum_{1\leqslant i\leqslant k}\mathbb{P}(J=i)
  = \mathbb{P}(T>t)\,\mathbb{P}(J\leqslant k).
\]
Tomando o complementar de $\{T>t\}$ em relação ao evento $\{J\leqslant k\}$, concluímos que
\[ 
  F_{T,J}(t,k)=F_T(t)F_J(k), \qquad t\geqslant 0 \ \  \text{e}\ \  k\in\mathbb{R}.
\] 
Para $t<0$, ambos os lados são zero.
Pela definição de independência, $T$ e $J$ são, portanto, independentes.
Em particular, saber qual relógio tocou primeiro
não altera a distribuição do tempo até o toque, 
isto é, $\mathbb{P}(T\leqslant t\mid J=i) =\mathbb{P}(T\leqslant t)$.
Observe que \eqref{eq-aux2-ex-corrida-despertadores} diz que o tempo mínimo tem taxa total
$\Lambda$, e a identidade do vencedor, segundo \eqref{eq-aux3-ex-corrida-despertadores}, 
é sorteada com probabilidades proporcionais às taxas individuais.

\bigskip 

Podemos agora usar a fórmula do vencedor como um método de amostragem:
escolheremos as taxas de modo que o índice vencedor tenha a distribuição
\emph{softmax} de uma lista de escores. Consideremos escores (ou pontuações ou log-pesos) 
$s_1,\ldots,s_n\in\mathbb{R}$. Como as taxas na distribuição exponencial precisam ser positivas, 
tomamos $\lambda_i=e^{s_i}$. Já que a função exponencial é positiva e estritamente crescente, esta transformação 
fornece taxas e preserva a ordem dos escores. Observe que agora, os valores
$\lambda_i \equiv e^{s_i}$ podem ser vistos como  pesos não normalizados 
da distribuição de Gibbs definida pelos escores. 
Substituindo essas taxas em \eqref{eq-aux3-ex-corrida-despertadores},
obtemos
\[
  \mathbb{P}(J=i)
  =\frac{\lambda_i}{\Lambda}
  =\frac{e^{s_i}}{\sum_{j=1}^n e^{s_j}}.
\]
Essa é a distribuição \emph{softmax} dos escores. Portanto, para sortear um
índice com essas probabilidades, basta gerar, independentemente,
$X_i\sim\mathrm{Exp}(e^{s_i})$ e escolher o índice do menor valor.

Essa distribuição conecta a corrida a dois exemplos vistos antes. Nas notas
sobre propriedades da probabilidade, ela aparece como a distribuição de Gibbs
obtida ao normalizar os pesos $e^{s_i}$. Nas notas sobre vetores discretos,
quando os escores do problema de atenção são $s_i$, os pesos são
$\alpha_i=e^{s_i}/\sum_{j=1}^n e^{s_j}$ e formam a combinação convexa
$\sum_i\alpha_i v_i$. A corrida usa esses mesmos pesos para sortear o índice
$J$, de modo que $\mathbb{P}(J=i)=\alpha_i$.
\end{example}

\begin{remark}[O truque Gumbel-Max]\label{rem:gumbel-max}
Há uma forma de fazer o mesmo sorteio usando apenas variáveis uniformes.
Para obter uma exponencial de taxa um, procuramos uma transformação cuja
probabilidade de exceder $t$ seja $e^{-t}$. Uma uniforme em $(0,1)$ fica
abaixo de $e^{-t}$ com essa probabilidade; como $-\log$ é decrescente, a
transformação $-\log U$ converte essa condição em $-\log U>t$.
Depois, reescalamos os tempos para obter as taxas $e^{s_i}$. Sejam
$U_1,\ldots,U_n:\Omega\to(0,1)$ variáveis aleatórias uniformes e
independentes, e defina
\[
  E_i=-\log U_i.
\]
Para $t\geqslant0$, a transformação fornece
\[
  \mathbb{P}(E_i>t)
  =\mathbb{P}(U_i<e^{-t})=e^{-t},
\]
de modo que $E_i\sim\mathrm{Exp}(1)$. Para verificar a independência, note
que, para $t_1,\ldots,t_n\geqslant0$,
\[
  \mathbb{P}(E_1\leqslant t_1,\ldots,E_n\leqslant t_n)
  =\mathbb{P}(U_1\geqslant e^{-t_1},\ldots,U_n\geqslant e^{-t_n})
  =\prod_{i=1}^n(1-e^{-t_i}).
\]
A segunda igualdade usa a independência das uniformes. Se algum $t_i<0$,
tanto a probabilidade conjunta quanto o produto das marginais são zero.
Assim, a função de distribuição conjunta se fatora, e $E_1,\ldots,E_n$
são independentes. Ponha
\[
  T_i=\frac{E_i}{e^{s_i}}.
\]
Para $t\geqslant0$, temos
$\mathbb{P}(T_i>t)=\mathbb{P}(E_i>e^{s_i}t)=e^{-e^{s_i}t}$.
Logo $T_i$ tem distribuição exponencial de taxa $e^{s_i}$. A fatoração
da função de distribuição conjunta continua válida após esse
reescalonamento de cada limiar, de modo que os tempos permanecem
independentes. Pelo \Cref{ex:corrida-exponenciais-softmax}, o índice que
minimiza os $T_i$ tem a distribuição softmax.

Para escrever essa regra como uma maximização, aplicamos $-\log$, que é
estritamente decrescente nos números positivos. Portanto,
\[
  \operatorname*{argmin}_{1\leqslant i\leqslant n}T_i
  =\operatorname*{argmax}_{1\leqslant i\leqslant n}(-\log T_i)
  =\operatorname*{argmax}_{1\leqslant i\leqslant n}(s_i+G_i),
\]
onde $G_i=-\log(-\log U_i)$. Aqui, $\operatorname*{argmin}$ e
$\operatorname*{argmax}$ denotam os índices em que o mínimo e o máximo
são atingidos. Como o empate tem probabilidade zero, esses índices são
únicos com probabilidade $1$. Os ruídos $G_i$ são independentes, pois
suas transformações são estritamente crescentes em cada $U_i$, e a
fatoração das funções de distribuição é preservada.
Calculamos a distribuição de cada $G_i$ diretamente. Para $g\in\mathbb{R}$,
\begin{align*}
  F_{G_i}(g)
  &=\mathbb{P}\bigl(-\log(-\log U_i)\leqslant g\bigr)
  \\[0.2cm]
  &=\mathbb{P}\bigl(U_i\leqslant e^{-e^{-g}}\bigr)
  =e^{-e^{-g}}.
\end{align*}
Essa é a distribuição Gumbel padrão. Assim, gerar ruídos Gumbel independentes,
somá-los aos escores e escolher o índice do maior valor é um algoritmo exato
para sortear uma posição com probabilidades softmax. Esse procedimento é
conhecido como \emph{truque Gumbel-Max}.\footnote{Veja C.~J.~Maddison,
\href{https://arxiv.org/abs/1602.05986}{\emph{A Poisson process model for
Monte Carlo}} (2016), para essa ligação entre corridas exponenciais
e o truque Gumbel-Max.}
\end{remark}

\subsection{Informação Conjunta e o Problema de Classificação}\label{subsec:classificacao-informacao-conjunta}

\begin{example}[Duas características e uma informação conjunta]\label{ex:classes-mesmas-marginais}
Queremos construir um exemplo em que cada característica, observada sozinha,
não ajude a identificar a classe, mas as duas características em conjunto
determinem a classe. Isso nos permitirá ver concretamente que as distribuições
marginais podem perder a informação sobre a interação entre variáveis.

Seja $(\Omega,\mathcal{F},\mathbb{P})$ um espaço de probabilidade e seja
$C:\Omega\to\{0,1\}$ uma variável aleatória com distribuição de Bernoulli
de parâmetro $1/2$. A variável $C$ divide o espaço amostral em dois grupos,
representados pelos eventos
\[
  \{C=0\}=\{\omega\in\Omega:C(\omega)=0\},
  \qquad
  \{C=1\}=\{\omega\in\Omega:C(\omega)=1\}.
\]
Esses eventos são disjuntos, sua união é $\Omega$, e
\[
  \mathbb{P}(C=0)=\mathbb{P}(C=1)=\tfrac12.
\]
Sejam ainda $X,Y:\Omega\to\mathbb{R}$ duas variáveis aleatórias. Para cada
resultado $\omega$, o par $(X(\omega),Y(\omega))$ registra duas características
desse resultado, enquanto $C(\omega)$ indica seu grupo. Um dos nossos objetivos é encontrar
a distribuição condicional conjunta dessas características (ou do vetor $(X,Y)$) 
dado um determinado grupo.
Este tipo de problema é muito comum quando estamos trabalhando em tarefas de classificação, 
já que, em geral, conhecemos algumas características do objeto que queremos
classificar, neste nosso exemplo seriam as variáveis aleatórias $X$ e $Y$, 
e queremos, com base nesses valores, inferir a que classe, indicada pela
variável aleatória $C$, nosso objeto pertence.

Vamos assumir que 
\[ 
\mathbb{P}\big( (X,Y)\in (0,1)\times (0,1)\big) =1.
\] 
Vamos considerar as seguintes regiões do quadrado aberto
$Q=(0,1)\times(0,1)$:
\begin{align*}
  R_0&=\bigl\{(x,y)\in Q :(x-\tfrac12)(y-\tfrac12)>0\bigr\},
  \\[0.3cm]
  R_1&=\bigl\{(x,y)\in Q :(x-\tfrac12)(y-\tfrac12)<0\bigr\}.
\end{align*}
\begin{figure}[H]
\centering
\begin{tikzpicture}[
  x=3.2cm,y=3.2cm,>=Stealth,
  eixo/.style={->,line width=0.6pt},
  divisoria/.style={draw=muted,densely dashed,line width=0.7pt},
  rotulo/.style={text=teal!80!black,font=\large}
]
  \foreach \indice/\deslocamento in {0/0,1/7.1} {
    \begin{scope}[xshift=\deslocamento cm]
      % Os quadrantes destacados dependem do sinal do produto.
      \ifnum\indice=0
        \fill[teal!20] (0,0) rectangle (0.5,0.5);
        \fill[teal!20] (0.5,0.5) rectangle (1,1);
        \node[rotulo] at (0.25,0.25) {$R_0$};
        \node[rotulo] at (0.75,0.75) {$R_0$};
      \else
        \fill[teal!20] (0,0.5) rectangle (0.5,1);
        \fill[teal!20] (0.5,0) rectangle (1,0.5);
        \node[rotulo] at (0.25,0.75) {$R_1$};
        \node[rotulo] at (0.75,0.25) {$R_1$};
      \fi

      \draw[draw=muted,line width=0.6pt] (0,0) rectangle (1,1);
      % O fundo branco separa visualmente os quadrantes abertos.
      \draw[preaction={draw=white,line width=2.2pt},divisoria]
        (0.5,0) -- (0.5,1);
      \draw[preaction={draw=white,line width=2.2pt},divisoria]
        (0,0.5) -- (1,0.5);
      \draw[eixo] (-0.08,0) -- (1.15,0) node[right] {$x$};
      \draw[eixo] (0,-0.08) -- (0,1.15) node[above] {$y$};

      \node[below left] at (0,0) {$0$};
      \foreach \valor/\texto in {0.5/{\tfrac12},1/{1}} {
        \draw (\valor,0.015) -- (\valor,-0.015);
        \draw (0.015,\valor) -- (-0.015,\valor);
        \node[below] at (\valor,-0.02) {$\texto$};
        \node[left] at (-0.02,\valor) {$\texto$};
      }
    \end{scope}
  }
\end{tikzpicture}

\caption{As regiões $R_0$ e $R_1$ são a união dos dois quadrantes sombreados.
As linhas divisórias $x=1/2$ e $y=1/2$, representadas por traços, e a borda
do quadrado estão excluídas das regiões.}
\label{fig:regioes-classificacao}
\end{figure}

É importante observar que os eventos $\{C=0\}$ e $\{C=1\}$ são subconjuntos de $\Omega$, 
enquanto as regiões $R_0$ e $R_1$ são subconjuntos de $\mathbb{R}^2$. 
Vamos relacioná-los abaixo no momento de construir a distribuição
condicional. Isto será feito de modo que, dado $\{C=0\}$, o par $(X,Y)$ pertença a $R_0$ quase certamente,
com respeito à distribuição condicional conjunta. 
Analogamente dado $\{C=1\}$ o par $(X,Y)$ tomará valores em $R_1$, quase certamente,
também com respeito à distribuição condicional conjunta.  

Para descrever a distribuição das características determinadas pelo vetor aleatório $(X,Y)$ 
em cada classe, usamos a medida condicional ao evento $\{C=c\}$, com $c=0,1$.
Para facilitar a notação vamos denotar estas medidas condicionais por $\mathbb{P}_{c}$.
Desta forma temos para todo $B\in\mathcal{F}$ que 
\[
  \mathbb{P}_c(B)=\mathbb{P}(B\mid C=c)
  =
  \frac{\mathbb{P}(B\cap\{C=c\})}{\mathbb{P}(C=c)}
  =
  \frac{\mathbb{P}(B\cap\{C=c\})}{1/2}
  =
  2\,\mathbb{P}(B\cap\{C=c\}).
\]
Vamos supor que sob $\mathbb{P}_c$ o vetor aleatório
$(X,Y)$ tenha a densidade conjunta $h_c:\mathbb{R}^2 \to\mathbb{R}$, onde
para cada $(x,y)\in\mathbb{R}^2$ temos
\[
  h_0(x,y)=2\cdot\mathds{1}_{R_0}(x,y),
  \qquad
  h_1(x,y)=2\cdot \mathds{1}_{R_1}(x,y).
\]
Para verificar que as funções $h_0$ e $h_1$ são de fato funções de densidade basta
observar que ambas são Borel-mensuráveis, não negativas e satisfazem 
\[
  \iint_{\mathbb{R}^2}h_c(x,y)\,dx\,dy
  =2\,\text{área}(R_c)=2\cdot\tfrac12=1,
  \qquad c\in\{0,1\}.
\]
Portanto, cada uma define uma distribuição de probabilidade para $(X,Y)$
dentro da classe correspondente.
A distribuição de $C$ e essas duas distribuições condicionais determinam a distribuição 
conjunta de $(C,X,Y)$. Mais precisamente, para todo boreliano
$S\subseteq\mathbb{R}^2$,
\[
  \mathbb{P}\bigl(C=c,(X,Y)\in S\bigr)
  =\tfrac12\,\mathbb{P}_c\bigl((X,Y)\in S\bigr)
  =\tfrac12\iint_S h_c(x,y)\,dx\,dy.
\]


Para determinar a distribuição de $X$ em cada classe, calculamos sua
densidade marginal sob $\mathbb{P}_c$. Pela fórmula
\eqref{eq:densidade-marginal-x}, uma versão dessa densidade é
\[
  h_{X,c}(x)=\int_{\mathbb{R}}h_c(x,y)\,dy,
  \qquad c\in\{0,1\}.
\]

Fixemos $x\in(0,1/2)$. Neste caso, $x-1/2<0$, portanto pelas definições de $R_0$ e $R_1$
temos as seguintes seções verticais:
\[
  \{y\in\mathbb{R}:(x,y)\in R_0\}=\bigl(0,\tfrac{1}{2}\bigr)
  \qquad \text{e}\qquad 
  \{y\in\mathbb{R}:(x,y)\in R_1\}=\bigl(\tfrac{1}{2},1\bigr).
\]
Como $h_c=2\cdot\mathds{1}_{R_c}$, segue que, para todo $y\in\mathbb{R}$,
\[
  h_0(x,y)=2\cdot\mathds{1}_{\left(0\,,\,\tfrac{1}{2}\right)}(y),
  \qquad
  h_1(x,y)=2\cdot\mathds{1}_{\left( \tfrac{1}{2}\,,\, 1\right)}(y).
\]
Essas seções estão representadas em vermelho na
\Cref{fig:secoes-verticais-classificacao}.
\begin{figure}[H]
\centering
\begin{tikzpicture}[
  x=3.6cm,y=3.6cm,>=Stealth,
  eixo/.style={->,thick},
  divisoria/.style={dashed,muted},
  secao/.style={red!75,thick},
  extremo/.style={draw=red!75,fill=white,line width=0.8pt}
]
  % O valor desenhado de x_0 é 0.3, pertencente a (0,1/2).
  \foreach \indice/\deslocamento in {0/0,1/7.1} {
    \begin{scope}[xshift=\deslocamento cm]
      \ifnum\indice=0
        \fill[teal!15] (0,0) rectangle (0.5,0.5);
        \fill[teal!15] (0.5,0.5) rectangle (1,1);
        \node at (0.15,0.25) {$R_0$};
        \node at (0.75,0.75) {$R_0$};
      \else
        \fill[teal!15] (0,0.5) rectangle (0.5,1);
        \fill[teal!15] (0.5,0) rectangle (1,0.5);
        \node at (0.15,0.75) {$R_1$};
        \node at (0.75,0.25) {$R_1$};
      \fi

      \draw[muted] (0,0) rectangle (1,1);
      \draw[divisoria] (0.5,0) -- (0.5,1);
      \draw[divisoria] (0,0.5) -- (1,0.5);
      \draw[eixo] (-0.08,0) -- (1.15,0) node[right] {$x$};
      \draw[eixo] (0,-0.08) -- (0,1.15) node[above] {$y$};

      % A linha tracejada fixa x=x_0; a parte contínua é a seção de R_c.
      \draw[red!75,dashed] (0.3,0) -- (0.3,1.07);
      \draw[red!75,dashed] (0,0.5) -- (0.3,0.5);
      \ifnum\indice=0
        \draw[secao] (0.3,0) -- (0.3,0.5);
        \draw[extremo] (0.3,0) circle (1.8pt);
        \draw[extremo] (0.3,0.5) circle (1.8pt);
        \node at (0.5,1.3) {$R_0$};
      \else
        \draw[red!75,dashed] (0,1) -- (0.3,1);
        \draw[secao] (0.3,0.5) -- (0.3,1);
        \draw[extremo] (0.3,0.5) circle (1.8pt);
        \draw[extremo] (0.3,1) circle (1.8pt);
        \node at (0.5,1.3) {$R_1$};
      \fi

      \node[below left=2pt] at (0,0) {$0$};
      \node[below=3pt] at (0.3,0) {$x_0$};
      \foreach \valor/\texto in {0.5/{\tfrac12},1/{1}} {
        \node[below=3pt] at (\valor,0) {$\texto$};
        \node[left=3pt] at (0,\valor) {$\texto$};
      }
    \end{scope}
  }
\end{tikzpicture}

\captionsetup{width=\textwidth}
\caption{Seções verticais de $R_0$ e $R_1$ para um valor fixado
$x_0\in(0,1/2)$. Os segmentos vermelhos correspondem aos intervalos
$0<y<1/2$ e $1/2<y<1$, respectivamente. Os círculos abertos indicam os
extremos excluídos.}
\label{fig:secoes-verticais-classificacao}
\end{figure}
Consequentemente,
\begin{align}
  h_{X,0}(x)&=\int_{\mathbb{R}}h_0(x,y)\,dy
  =\int_0^{1/2}2\,dy=1, \qquad \forall x\in \left(0,\tfrac{1}{2}\right)
  \label{eq:marginal-x-classe0-esquerda}
  \\[0.2cm]
  h_{X,1}(x)&=\int_{\mathbb{R}}h_1(x,y)\,dy
  =\int_{1/2}^{1}2\,dy=1, \qquad \forall x\in \left(0,\tfrac{1}{2}\right).
  \label{eq:marginal-x-classe1-esquerda}
\end{align}

Fixemos agora $x\in(1/2,1)$. Agora, $x-1/2>0$, portanto pelas definições
de $R_0$ e $R_1$ temos as seguintes seções verticais:
\[
  \{y\in\mathbb{R}:(x,y)\in R_0\}=\bigl(\tfrac{1}{2},1\bigr)
  \qquad \text{e}\qquad
  \{y\in\mathbb{R}:(x,y)\in R_1\}=\bigl(0,\tfrac{1}{2}\bigr).
\]
Como $h_c=2\cdot\mathds{1}_{R_c}$, segue que, para todo $y\in\mathbb{R}$,
\[
  h_0(x,y)=2\cdot\mathds{1}_{\left(\tfrac{1}{2}\,,\,1\right)}(y),
  \qquad
  h_1(x,y)=2\cdot\mathds{1}_{\left(0\,,\,\tfrac{1}{2}\right)}(y).
\]
Essas seções estão representadas em vermelho na
\Cref{fig:secoes-verticais-classificacao-direita}.
\begin{figure}[H]
\centering
\begin{tikzpicture}[
  x=3.6cm,y=3.6cm,>=Stealth,
  eixo/.style={->,thick},
  divisoria/.style={dashed,muted},
  secao/.style={red!75,thick},
  extremo/.style={draw=red!75,fill=white,line width=0.8pt}
]
  % O valor desenhado de x_0 é 0.7, pertencente a (1/2,1).
  \foreach \indice/\deslocamento in {0/0,1/7.1} {
    \begin{scope}[xshift=\deslocamento cm]
      \ifnum\indice=0
        \fill[teal!15] (0,0) rectangle (0.5,0.5);
        \fill[teal!15] (0.5,0.5) rectangle (1,1);
        \node at (0.25,0.25) {$R_0$};
        \node at (0.9,0.75) {$R_0$};
      \else
        \fill[teal!15] (0,0.5) rectangle (0.5,1);
        \fill[teal!15] (0.5,0) rectangle (1,0.5);
        \node at (0.25,0.75) {$R_1$};
        \node at (0.9,0.25) {$R_1$};
      \fi

      \draw[muted] (0,0) rectangle (1,1);
      \draw[divisoria] (0.5,0) -- (0.5,1);
      \draw[divisoria] (0,0.5) -- (1,0.5);
      \draw[eixo] (-0.08,0) -- (1.15,0) node[right] {$x$};
      \draw[eixo] (0,-0.08) -- (0,1.15) node[above] {$y$};

      % A linha tracejada fixa x=x_0; a parte contínua é a seção de R_c.
      \draw[red!75,dashed] (0.7,0) -- (0.7,1.07);
      \draw[red!75,dashed] (0,0.5) -- (0.7,0.5);
      \ifnum\indice=0
        \draw[red!75,dashed] (0,1) -- (0.7,1);
        \draw[secao] (0.7,0.5) -- (0.7,1);
        \draw[extremo] (0.7,0.5) circle (1.8pt);
        \draw[extremo] (0.7,1) circle (1.8pt);
        \node at (0.5,1.3) {$R_0$};
      \else
        \draw[secao] (0.7,0) -- (0.7,0.5);
        \draw[extremo] (0.7,0) circle (1.8pt);
        \draw[extremo] (0.7,0.5) circle (1.8pt);
        \node at (0.5,1.3) {$R_1$};
      \fi

      \node[below left=2pt] at (0,0) {$0$};
      \node[below=3pt] at (0.7,0) {$x_0$};
      \foreach \valor/\texto in {0.5/{\tfrac12},1/{1}} {
        \node[below=3pt] at (\valor,0) {$\texto$};
        \node[left=3pt] at (0,\valor) {$\texto$};
      }
    \end{scope}
  }
\end{tikzpicture}

\captionsetup{width=\textwidth}
\caption{Seções verticais de $R_0$ e $R_1$ para um valor fixado
$x_0\in(1/2,1)$. Os segmentos vermelhos correspondem aos intervalos
$1/2<y<1$ e $0<y<1/2$, respectivamente. Os círculos abertos indicam os
extremos excluídos.}
\label{fig:secoes-verticais-classificacao-direita}
\end{figure}
Consequentemente,
\begin{align}
  h_{X,0}(x)&=\int_{\mathbb{R}}h_0(x,y)\,dy
  =\int_{1/2}^{1}2\,dy=1, \qquad \forall x\in\left(\tfrac{1}{2},1\right)
  \label{eq:marginal-x-classe0-direita}
  \\[0.2cm]
  h_{X,1}(x)&=\int_{\mathbb{R}}h_1(x,y)\,dy
  =\int_0^{1/2}2\,dy=1, \qquad \forall x\in\left(\tfrac{1}{2},1\right).
  \label{eq:marginal-x-classe1-direita}
\end{align}

Portanto, pelas igualdades \eqref{eq:marginal-x-classe0-esquerda},
\eqref{eq:marginal-x-classe1-esquerda}, \eqref{eq:marginal-x-classe0-direita}
e \eqref{eq:marginal-x-classe1-direita}, para cada $c\in\{0,1\}$ fixado e
cada $x\in(0,1)\setminus\{\tfrac{1}{2}\}$ temos
\[
  \int_{\mathbb{R}}h_c(x,y)\,dy=2\cdot\tfrac{1}{2}=1.
\]
Para $x\notin(0,1)$, essa integral é zero. Para $x=1/2$, ela também é
zero, pois fixamos as densidades conjuntas como zero nas linhas divisórias.
Podemos, porém, escolher outra versão da densidade marginal, atribuindo-lhe
o valor $1$ nesse único ponto, sem alterar a distribuição de $X$.
O mesmo argumento vale para os cortes horizontais, já que
$h_c(x,y)=h_c(y,x)$. Assim, passaremos a usar as versões marginais
\[
  h_{X,c}(x)=\mathds{1}_{(0,1)}(x),
  \qquad
  h_{Y,c}(y)=\mathds{1}_{(0,1)}(y),
  \qquad c\in\{0,1\}.
\]
Essas versões coincidem com as integrais marginais em quase todo ponto.
Em particular, sob ambas as medidas $\mathbb{P}_0$ e $\mathbb{P}_1$,
as duas variáveis aleatórias $X$ e $Y$ têm, individualmente, distribuição uniforme em $(0,1)$.

Desta forma temos
\[
\mathbb{P}_0(X\in D)=\mathbb{P}_1(X\in D), \qquad \forall D\in\mathscr{B}(\mathbb{R}).
\] 
A mesma igualdade também é válida para a variável aleatória $Y$. 

Por exemplo, saber que $X$ pertence a um intervalo $I\subseteq(0,1)$ 
de comprimento positivo mantém as probabilidades das duas classes 
iguais às probabilidades iniciais:
\[
  \mathbb{P}(C=0\mid X\in I)
  =\mathbb{P}(C=1\mid X\in I)=\tfrac12.
\]
Para justificar essa igualdade, denotemos por $\ell(I)>0$ o comprimento de
$I$. Como a distribuição de $X$ sob $\mathbb{P}_c$ é uniforme em $(0,1)$,
para cada $c\in\{0,1\}$ temos
\[
  \mathbb{P}(X\in I\mid C=c)
  =\mathbb{P}_c(X\in I)
  =\int_I h_{X,c}(x)\,dx
  =\int_I 1\,dx=\ell(I).
\]
Os eventos $\{C=0\}$ e $\{C=1\}$ formam uma partição de $\Omega$.
Logo, pela fórmula da probabilidade total,
\begin{align*}
  \mathbb{P}(X\in I)
  &=\sum_{c=0}^{1}\mathbb{P}(X\in I\mid C=c)\,\mathbb{P}(C=c)
  \\[0.2cm]
  &=\ell(I)\cdot\tfrac12+\ell(I)\cdot\tfrac12
  =\ell(I)>0.
\end{align*}
Portanto, podemos condicionar ao evento $\{X\in I\}$. Pela definição de
probabilidade condicional e pela fórmula de Bayes, para cada
$c\in\{0,1\}$ obtemos
\begin{align*}
  \mathbb{P}(C=c\mid X\in I)
  &=\frac{\mathbb{P}(\{C=c\}\cap\{X\in I\})}{\mathbb{P}(X\in I)}
  \\[0.2cm]
  &=\frac{\mathbb{P}(X\in I\mid C=c)\,\mathbb{P}(C=c)}
    {\mathbb{P}(X\in I)}
  \\[0.2cm]
  &=\frac{\ell(I)\cdot\tfrac12}{\ell(I)}=\tfrac12.
\end{align*}
O cálculo mostra por que as probabilidades das classes permanecem iguais:
o evento $\{X\in I\}$ tem a mesma probabilidade condicional em ambas as
classes.
O mesmo vale quando a informação é sobre $Y$. Assim, informações sobre uma
única coordenada não alteram as probabilidades dos grupos; para recuperar
o valor de $C$, precisamos considerar como as duas coordenadas se combinam.

Para quantificar a limitação descrita acima, podemos usar uma
regra de previsão baseada exclusivamente na variável aleatória $X$. 
Uma regra de previsão é simplesmente uma função Borel-mensurável 
$g:\mathbb{R}\to\{0,1\}$ e funciona da seguinte maneira. 
Ao observar o valor $X(\omega)=x$, a classe prevista será a classe $g(x)$. 
Uma regra de previsão às vezes é chamada de uma regra ``\emph{determinística}'', 
porque, para um mesmo valor $x$, a função $g$ sempre produz a mesma previsão. 

Denotemos por $\widehat{C}$ a previsão obtida a partir de $X$:
\[
  \widehat{C}:\Omega\to\{0,1\},
  \qquad \widehat{C}(\omega)=g(X(\omega)).
\]
Como $X$ e $g$ são mensuráveis, $\widehat{C}$ é uma variável aleatória.
Observe que $C$ indica a classe verdadeira, enquanto $\widehat{C}$ indica
a classe prevista pela regra ``determinística''.

Já que $g$ é mensurável e assume apenas dois valores então o boreliano 
\[
  D=g^{-1}(\{0\})=\{x\in\mathbb{R}:g(x)=0\},
\]
determina $g$ e vice-versa já que 
\[
  g(x)=
  \begin{cases}
    0, & x\in D,\\[0.1cm]
    1, & x\notin D.
  \end{cases}
\]
Como $X$ é uma variável aleatória e $D$ é boreliano, temos imediatamente que
\[
  \{\widehat{C}=0\}=\{X\in D\}\in\mathcal{F}
	\qquad \text{e}\qquad 
  \{\widehat{C}=1\}=\{X\notin D\}\in\mathcal{F}.
\]

A regra ``determinística'' acerta quando $\widehat{C}(\omega)=C(\omega)$. Portanto, o evento de
acerto é a união disjunta
\[
  \{\widehat{C}=C\}
  =\bigl(\{C=0\}\cap\{X\in D\}\bigr)
   \cup\bigl(\{C=1\}\cap\{X\notin D\}\bigr).
\]
Usando a definição de $\mathbb{P}_c$ e a igualdade das distribuições
marginais nas duas classes, obtemos
\begin{align*}
  \mathbb{P}(\widehat{C}=C)
  &=\mathbb{P}(C=0,X\in D)+\mathbb{P}(C=1,X\notin D)
  \\[0.2cm]
  &=\tfrac12\,\mathbb{P}_0(X\in D)
    +\tfrac12\,\mathbb{P}_1(X\notin D)
  \\[0.2cm]
  &=\tfrac12\,\mathbb{P}_0(X\in D)
    +\tfrac12\bigl(1-\mathbb{P}_0(X\in D)\bigr)
  =\tfrac12,
\end{align*}
onde na penúltima igualdade usamos que 
$\mathbb{P}_1(X\notin D)=\mathbb{P}_0(X\notin D)=1-\mathbb{P}_0(X\in D)$.

Logo, toda regra determinística mensurável baseada somente na variável aleatória $X$ acerta com
probabilidade $1/2$.
O mesmo cálculo vale para qualquer regra determinística baseada somente em
$Y$. 


\bigskip 

Por outro lado, vamos mostrar que é possível construir uma regra determinística
mensurável baseada no vetor $(X,Y)$ que classifica corretamente com
probabilidade $1$. Agora, uma regra de previsão é uma função
$G:\mathbb{R}^2\to\{0,1\}$: ao observar
$(X(\omega),Y(\omega))=(x,y)$, a classe prevista será $G(x,y)$.

Para escolher essa função, usamos as definições de $R_0$ e $R_1$. Em $R_0$,
as duas coordenadas são simultaneamente menores que $1/2$ ou simultaneamente
maiores que $1/2$, de modo que $(x-\tfrac12)(y-\tfrac12)>0$. Em $R_1$,
uma coordenada é menor que $1/2$ e a outra é maior, de modo que esse produto
é negativo. Assim, definimos, para todo $(x,y)\in\mathbb{R}^2$,
\[
  G(x,y)=
  \begin{cases}
    0, & (x-\tfrac12)(y-\tfrac12)\geqslant 0,\\[0.2cm]
    1, & (x-\tfrac12)(y-\tfrac12)<0.
  \end{cases}
\]
Essa definição atribui a previsão $0$ às linhas divisórias $x=1/2$ e
$y=1/2$. Como a função $(x,y)\mapsto(x-\tfrac12)(y-\tfrac12)$ é contínua,
o conjunto $G^{-1}(\{0\})$ é fechado e o conjunto $G^{-1}(\{1\})$ é aberto.
Ambos são borelianos de $\mathbb{R}^2$, logo $G$ é Borel-mensurável.

Denotemos por $\widehat{C}_{X,Y}$ a previsão obtida a partir das duas
características:
\[
  \widehat{C}_{X,Y}:\Omega\to\{0,1\},
  \qquad \widehat{C}_{X,Y}(\omega)=G(X(\omega),Y(\omega)).
\]
Como $(X,Y)$ e $G$ são mensuráveis, $\widehat{C}_{X,Y}$ é uma variável
aleatória. Para cada $c\in\{0,1\}$, o evento em que a regra prevê a classe
$c$ é
\[
  \{\widehat{C}_{X,Y}=c\}
  =\{(X,Y)\in G^{-1}(\{c\})\}\in\mathcal{F}.
\]

Pela construção de $G$, temos $R_c\subseteq G^{-1}(\{c\})$ para
$c\in\{0,1\}$. Além disso, a densidade $h_c$ é nula fora de $R_c$.
Portanto, sob a medida condicional $\mathbb{P}_c$, a regra prevê a classe
$c$ com probabilidade $1$:
\begin{align*}
  \mathbb{P}_c(\widehat{C}_{X,Y}=c)
  &=\iint_{G^{-1}(\{c\})}h_c(x,y)\,dx\,dy
  \\[0.2cm]
  &=\iint_{R_c}h_c(x,y)\,dx\,dy
  =2\cdot\tfrac12=1.
\end{align*}
Para obter a probabilidade de acerto sob $\mathbb{P}$, escrevemos novamente
o evento de acerto como uma união disjunta:
\[
  \{\widehat{C}_{X,Y}=C\}
  =\bigl(\{C=0\}\cap\{\widehat{C}_{X,Y}=0\}\bigr)
   \cup\bigl(\{C=1\}\cap\{\widehat{C}_{X,Y}=1\}\bigr).
\]
Usando a definição das medidas condicionais, obtemos
\begin{align*}
  \mathbb{P}(\widehat{C}_{X,Y}=C)
  &=\mathbb{P}(C=0,\widehat{C}_{X,Y}=0)
    +\mathbb{P}(C=1,\widehat{C}_{X,Y}=1)
  \\[0.2cm]
  &=\tfrac12\,\mathbb{P}_0(\widehat{C}_{X,Y}=0)
    +\tfrac12\,\mathbb{P}_1(\widehat{C}_{X,Y}=1)
  \\[0.2cm]
  &=\tfrac12\cdot1+\tfrac12\cdot1=1.
\end{align*}
Como as linhas divisórias têm área zero e a distribuição conjunta de
$(X,Y)$ sob cada $\mathbb{P}_c$ tem densidade, o evento
$\{X=1/2\}\cup\{Y=1/2\}$ tem probabilidade zero sob cada $\mathbb{P}_c$.
Pela fórmula da probabilidade total, esse evento também tem probabilidade
zero sob $\mathbb{P}$. Portanto, qualquer escolha da previsão nessas linhas
produz a mesma probabilidade de acerto.
Assim, com probabilidade $1$, o valor de $C$ pode ser recuperado a partir
do par $(X,Y)$.
A interação entre as coordenadas contém toda a informação de classe, embora
cada marginal, separadamente, não contenha nenhuma.

Essa construção também mostra o que se perde ao substituir, dentro de cada
classe, a densidade conjunta pelo produto das densidades marginais. Em ambas
as classes esse produto é
$h_{X,c}(x)h_{Y,c}(y)=\mathds{1}_{Q}(x,y)$, uma densidade uniforme em $Q$.
Como esse produto não depende de $c$, a substituição atribui a mesma
distribuição conjunta às duas classes e elimina a informação que permitia
recuperar $C$ a partir do vetor $(X,Y)$.

\end{example}

\subsection{Dimensão e escassez de amostras locais}\label{subsec:dimensao-amostras-locais}

\begin{example}[Dimensão e escassez de amostras locais]\label{ex:dimensao-amostras-locais}
  No exemplo anterior, usamos o vetor $(X,Y)$ para representar duas
  características de um objeto. Agora, fixemos um inteiro $d\geqslant1$ e
  consideremos objetos descritos por $d$ características reais. Cada objeto
  será representado por um vetor $x=(x_1,\ldots,x_d)\in\mathbb{R}^d$, no qual
  a coordenada $x_i$ registra o valor da $i$-ésima característica. Assim, o
  número de características é a dimensão do espaço em que representamos os
  objetos.

  Em algumas tarefas de inteligência artificial, procuramos, entre os dados
  disponíveis, objetos com características próximas às de um objeto de
  referência. Neste contexto, chamaremos de \emph{amostra} cada vetor de dados
  disponível e de \emph{consulta} o vetor fixo
  $q=(q_1,\ldots,q_d)\in\mathbb{R}^d$ que representa o objeto de referência.
  Escolhida uma tolerância $\varepsilon>0$, diremos que uma amostra $x$ está
  próxima da consulta $q$ em todas as características quando
  \[
    |x_i-q_i|<\varepsilon,\qquad\forall i\in\{1,\ldots,d\}.
  \]
  Queremos estudar como a dimensão $d$ afeta a probabilidade de obter uma
  amostra que satisfaça simultaneamente essas $d$ condições.

  Para isso, vamos usar um modelo probabilístico simples. Seja
  $(\Omega,\mathcal{F},\mathbb{P})$ um espaço de probabilidade e sejam
  $X_1,\ldots,X_d:\Omega\to\mathbb{R}$ variáveis aleatórias independentes,
  cada uma com distribuição uniforme em $(0,1)$, isto é, $X_i\sim U(0,1)$. O vetor aleatório
  \[
    \mathbf{X}:\Omega\to\mathbb{R}^d,
    \qquad \mathbf{X}(\omega)=(X_1(\omega),\ldots,X_d(\omega)),
  \]
  modela a geração de uma amostra: para cada resultado $\omega$, o vetor
  $\mathbf{X}(\omega)$ é a amostra obtida, e $X_i(\omega)$ é o valor de sua
  $i$-ésima característica. Esse vetor pertence ao cubo $(0,1)^d$ quase
  certamente.

  Neste exemplo, fixaremos a consulta no centro do cubo,
  $q=(1/2,\ldots,1/2)$, e usaremos a tolerância $\varepsilon=1/4$.
  A região dos vetores próximos dessa consulta é
  \[
    V_d=\prod_{i=1}^{d}\left(q_i-\frac14,q_i+\frac14\right)
    =\left(\frac14,\frac34\right)^d.
  \]
  O evento de que a amostra gerada esteja próxima dessa consulta em todas
  as características é dado por
  \[
    A_d=\{\mathbf{X}\in V_d\}
    =\bigcap_{i=1}^{d}\left\{|X_i-q_i|<\frac14\right\}.
  \]
  Observe que $V_d$ é uma região de $\mathbb{R}^d$, enquanto $A_d$ é um
  evento de $\Omega$: ele reúne os resultados $\omega$ para os quais
  $\mathbf{X}(\omega)$ pertence a $V_d$.

  Para calcular $\mathbb{P}(A_d)$, começamos pela condição em uma única
  característica. Como $q_i=1/2$ e $X_i\sim U(0,1)$, para cada
  $i\in\{1,\ldots,d\}$ temos
  \[
    \mathbb{P}\left(|X_i-q_i|<\frac14\right)
    =\mathbb{P}\left(\frac14<X_i<\frac34\right)
    =\int_{1/4}^{3/4}1\,dx=\frac12.
  \]
  Pela independência das variáveis $X_1,\ldots,X_d$, os eventos
  $\{|X_i-q_i|<1/4\}$ são independentes. Portanto,
  \begin{equation}\label{eq:probabilidade-regiao-central}
    \mathbb{P}(A_d)
    =\prod_{i=1}^{d}\mathbb{P}\left(|X_i-q_i|<\frac14\right)
    =2^{-d}.
  \end{equation}
  Geometricamente, $V_d$ tem volume $2^{-d}$ e o cubo $(0,1)^d$ tem volume
  $1$. Como $\mathbf{X}$ tem distribuição uniforme nesse cubo,
  $\mathbb{P}(A_d)$ coincide com a fração do volume ocupada por $V_d$.

  Para estudar a chance de sucesso com uma coleção de amostras, fixe um
  inteiro $N\geqslant1$ e sejam
  $\mathbf{X}^{(k)}=(X_1^{(k)},\ldots,X_d^{(k)})$, para $k=1,\ldots,N$,
  vetores aleatórios definidos no mesmo espaço de probabilidade.
  Assumimos que as $Nd$ variáveis $X_i^{(k)}$ sejam todas independentes e
  uniformes em $(0,1)$. Aqui, o índice $i$ identifica a característica e o
  índice $k$ identifica a amostra.

  Para cada $k\in\{1,\ldots,N\}$, o evento de que a $k$-ésima amostra
  esteja próxima da consulta é
  \[
    A_d^{(k)}=\{\mathbf{X}^{(k)}\in V_d\}.
  \]
  Cada vetor $\mathbf{X}^{(k)}$ tem a mesma distribuição que $\mathbf{X}$.
  Logo, por \eqref{eq:probabilidade-regiao-central},
  $\mathbb{P}(A_d^{(k)})=2^{-d}$. Além disso, os eventos
  $A_d^{(1)},\ldots,A_d^{(N)}$ são independentes: cada um é determinado
  pelas coordenadas de um único vetor, e esses grupos de coordenadas são
  independentes entre si.

  Encontrar ao menos uma amostra próxima da consulta corresponde ao evento
  \[
    B_{d,N}=\bigcup_{k=1}^{N}A_d^{(k)}.
  \]
  Para calcular sua probabilidade, é conveniente considerar o evento
  complementar: nenhuma das $N$ amostras pertence a $V_d$. Como os
  complementares dos eventos $A_d^{(k)}$ também são independentes, temos
  \begin{align*}
    \mathbb{P}(B_{d,N}^{c})
    &=\mathbb{P}\left(\bigcap_{k=1}^{N}(A_d^{(k)})^c\right)
    \\[0.2cm]
    &=\prod_{k=1}^{N}\mathbb{P}((A_d^{(k)})^c)
    =(1-2^{-d})^N.
  \end{align*}
  Passando ao complementar e aplicando a desigualdade da união, obtemos
  \begin{equation}\label{eq:probabilidade-alguma-amostra-central}
    \begin{aligned}
      \mathbb{P}(B_{d,N})&=1-(1-2^{-d})^N
      \\[0.2cm]
      &\leqslant\sum_{k=1}^{N}\mathbb{P}(A_d^{(k)})=N2^{-d}.
    \end{aligned}
  \end{equation}

  Por exemplo, com $d=30$ características e $N=1\,000\,000$ amostras, a
  probabilidade de encontrar ao menos uma amostra próxima da consulta
  satisfaz
  \[
    \mathbb{P}(B_{30,1\,000\,000})
    \leqslant \frac{1\,000\,000}{2^{30}}
    < 0{,}001 = 0{,}1\%.
  \]
  Para uma dimensão $d$ fixada, se desejarmos
  $\mathbb{P}(B_{d,N})\geqslant1/2$, o número $N$ de amostras deve
  satisfazer, por \eqref{eq:probabilidade-alguma-amostra-central},
  \begin{equation}\label{eq:amostras-necessarias-regiao-central}
    \frac12\leqslant\mathbb{P}(B_{d,N})\leqslant N2^{-d},
    \qquad\text{e, portanto,}\qquad N\geqslant 2^{d-1}.
  \end{equation}
  Essa condição é necessária, pois foi obtida de uma cota superior. Para
  verificar se um valor de $N$ é suficiente, usamos a expressão exata
  $1-(1-2^{-d})^N$ em \eqref{eq:probabilidade-alguma-amostra-central}.

  Com a consulta e a tolerância escolhidas, cada característica adicional
  impõe uma condição independente de probabilidade $1/2$, reduzindo pela
  metade a probabilidade de uma única amostra estar próxima de $q$.
  A condição \eqref{eq:amostras-necessarias-regiao-central} impõe ao número
  de amostras uma cota inferior que cresce exponencialmente com a
  dimensão. Esse comportamento ilustra, no modelo considerado, um
  dos fenômenos conhecidos como \emph{maldição da dimensionalidade}.
\end{example}

\bibliographystyle{plain}
\bibliography{../arquivos-aux/referencias}
\end{document}
